\begin{savequote}
\includegraphics[width=3.5cm]{./images/wikipedia-pd/Battlefield_wounds_300px.jpg}

{\tiny\itshape Hans von Gersdorff, Feldbuch der Wundarznei, 1517.}
\end{savequote}
\chapter{[TRA] \emph{Spezielle Patientengruppe:} Trauma}
\label{chp:TRA}
% \AasRsN{RS~VII}

\emph{Maintainer: Michael Motal}

\minitoc


\Definition{\T{Trauma:} Als Trauma  bezeichnet man eine
\begin{inparaenum}
  \item Schädigung,
  \item Verletzung, oder
  \item Wunde,
\end{inparaenum}
die durch eine
Gewalteinwirkung physikalischer (mechanischer,  thermischer) oder chemischer
Art verursacht wurde.}

Die Untersuchung des traumatologischen Patienten
(Traumacheck) wird unter \prettyref{topic:traumacheck} erklärt. 

%Die Traumata lassen sich grob ind folgende Klassen Einteilen:
% 
% \begin{itemize}
%     \item Schädel-Hirn-Trauma (SHT)
%     \item Thorax trauma
%     \item Abdominal trauma / Bauchtrauma
%     \item Beckentrauma
%     \item Verletzungen der Wirbelsäule
%     \item Polytrauma
%     \item Verletzungen der Extremitäten
%     \item Verbrennungen
% \end{itemize}



%%%%%%%%%%%%%%%%%%%%%%%%%%%%%%%%%%%%%%%%%%%%%%%%%%%%%%%%%%%
%%%%%%%%%%%%%%%%%%%%%%%%%%%%%%%%%%%%%%%%%%%%%%%%%%%%%%%%%%%
%%%%%%%%%%%%%%%%%%%%%%%%%%%%%%%%%%%%%%%%%%%%%%%%%%%%%%%%%%%

\section{Einsatztaktische Überlegungen bei verunfallten Patienten}

\subsection{Unfallort}
Bevor die Patientenversorgung beginnen kann, müssen am Unfallort einige
Vorkehrungen für den weiteren Einsatzablauf getroffen werden. Die wichtigsten
Aufgaben sind:
\begin{itemize}
  \item \Es{Absichern} der Unfallstelle und Selbstschutz
  (Blaulicht, Warnblinkanlage, Uniformjacke anziehen, ev. Helm aufsetzen)
  \item Evtl. Nachalarmierung weiterer Einsatzkräfte:
  weitere Rettungsmittel bei mehreren Verletzten, Feuerwehr, Polizei
  (Alarmierung ausschließlich über das Journal, Meldung der Anzahl an Verletzten)
  \item Bei mehreren Verletzten Triage (siehe ...)
  \item Eingeklemmte Personen werden im Fahrzeug stabilisiert bis die Feuerwehr
  die Verunfallten rettet; Helme tragen!
  \item Crash-Rettung nur bei Fremdgefährdung, HWS stabilisieren!
\end{itemize}

\subsection{Unfallmechanismus}
Der Unfallmechanismus gibt wertvolle Hinweise auf wahrscheinliche bzw. mögliche
Verletzungen des Patienten. Der Rettungssanitäter kann durch eine systematische
Vorgehensweise bei der Patientenversorgung, die eine genaue Betrachtung des
Verletzungsmechanismus beinhaltet, mithelfen, die Schäden eines Traumas begrenzt
zu halten. Die wichtigsten Fragen dabei sind immer:
\begin{itemize}
    \item \emph{Was ist passiert?}
    \item \emph{Wie ist der Patient verletzt worden?}
\end{itemize}

\Layout{22}{Verkehrsunfälle}
{ Bei Verkehrsunfällen müssen bei der Einschätzung des Verletzungsbildes drei
Faktoren beachtet werden:
\begin{itemize}
  \item Deformierungsgrad des Fahrzeugs (Indiz für die involvierten Kräfte)
  \item Deformierung von Teilen der Fahrzeugkabine (Hinweis für Aufschlagpunkt
  des Körpers im Fahrzeug)
  \item Deformierung (Verletzungsmuster) des Patienten (Anzeichen dafür,
  welche Körperteile direkt aufgeprallt sind)
\end{itemize}
}{
\SymbolText
}{./images/teaser/doa.jpg}{Oft der Verlierer: Der
Fußgänger.}{Hauer.}



\subparagraph{Unfälle mit Fahrzeugen}

Für Unfälle mit Fahrzeugen gilt allgemein:

\begin{itemize}
  \item Die Geschwindigkeit des Fahrzeugs vor dem Unfall spielt eine wichtige 
  Rolle. Je schneller das Fahrzeug unterwegs  war, desto schwerer kann die
  Verletzung sein. \Z{Dabei steigt die Bewegungsenergie des Fahrzeugs  mit dem
  Quadrat der Geschwindigkeit an. Eine Verdoppelung der Geschwindigkeit  hat
  daher eine Vervierfachung der Energie zur Folge, welche im Falle eines
  Unfalls frei wird und die Verformung verursacht.}
  \item Die Gewalteinwirkung auf den Patienten entspricht in ihrer Richtung 
  i.\,d.\,R. der Einwirkung auf das Fahrzeug.
  \item Wird ein Insasse aus dem Fahrzeug geschleudert,
  steigert sich das Risiko schwerer Verletzungen (SHT,
  Wirbelsäule, Thorax, Abdomen etc.) um 300\,\%
  \cite{BoehmerTaschRett6}. Auch bei
  den im Auto verbliebenen Insassen ist mit schweren Verletzungen zu rechnen.
  \item Wenn ein Insasse getötet oder eingeklemmt wird, ist bei den weiteren 
  Insassen das Risiko schwerer Verletzungen erhöht.
  \item Knieverletzungen bei einem Autounfall können Hinweise auf Oberschenkel-
  und Beckenfrakturen sowie Hüftgelenksverrenkungen sein.
\end{itemize}

\subparagraph{Frontalzusammenstöße}

Bei Frontalzusammenstößen sollten folgende Merkmale
Alarmzeichen für schwere Verletzungen (Thoraxtrauma mit  Herzkontusion und
Pneumothorax, HWS-Trauma, Leber- und Milzruptur, Beckenfraktur) sein:
\begin{itemize}
    \item Knieanstoß am Armaturenbrett,
    \item Eingedrückte, leicht nach außen gewölbte Windschutzscheibe  (ev.
    Spinnennetzmuster) durch Anstoß mit dem Kopf,
    \item sichtbare Veränderungen des Radstandes (z. B. Rad zur Seite 
    geknickt, Achsenverschiebung um mehr als 30 cm am Unfallfahrzeug),
    \item Fahrzeugverformung um mehr als 50 cm.
\end{itemize}

\subparagraph{Seitenaufprall}

Beim Seitenaufprall sind typische schwere Verletzungen: HWS-Verrenkung, 
Thoraxtrauma, Verletzung der Aorta, seitenabhängig Leber- oder Milzruptur,  
Frakturen des Beckens. Bei einem Heckaufprall sollte man an HWS-Trauma 
(Schleudertrauma) denken und auch eine Frontalaufprallkomponente  (Insasse wird
nach vorne geworfen) mit einbeziehen (Lenkradkontusion).

\subparagraph{Fußgänger vs. Fahrzeug}

Wird ein \Es{Fußgänger von einem Fahrzeug (frontal) erfasst}
ist generell mit schweren Verletzungen (vor allem
Schädel-Hirn-Trauma, Wirbelsäule) zu rechnen. \Z{Zur
besseren Einschätzung kann man folgende Faustregel (nach
Böhmer et al. 2006) verwenden:
\begin{itemize}
\item bis 50\,\nicefrac{km}{h}: Aufschlagen des Kopfes auf die Kühlerhaube
\item bis 70\,\nicefrac{km}{h}: Aufschlagen des Kopfes auf die Windschutzscheibe
\item mehr als 70\,\nicefrac{km}{h}: der Fußgänger wird über das Dach geworfen
\end{itemize}}
Lebensbedrohliche Verletzungen sind bereits möglich, wenn
Radfahrer oder Fußgänger mit mehr als 30\,\nicefrac{km}{h} durch ein
Fahrzeug erfasst werden. Erkennbar kann dies durch
Bremsspuren, die bei trockener Fahrbahn länger als \VonBis{10}{20}\,m
sind, sein.

\Layout{2}{Sturz aus großer Höhe}{%
Bei Stürzen aus großen Höhen (z. B. Leiter, Gerüst) ist stets
an HWS-, Wirbelsäulen- und SHT zu denken. Bei Sturz auf die
gestreckten Beine findet man typischerweise eine Fraktur der
Fersenbeine und/oder eine Stauchungsfraktur der Wirbelsäule
(Th\,12/L\,1), sowie Schienbeinplateau-Frakturen. Ab einer
Fallhöhe von 6\,m ist von einer Polytraumatisierung
auszugehen.
}{
\begin{itemize}
  \item bla!
\end{itemize}
}{}{}{}


\Layout{60}{Stich- und Schussverletzungen}{%
Traumen, die durch spitze Gegenstände (Stich- und
Pfählungsverletzungen) oder Schusswaffenmunition
(Schussverletzungen) hervorgerufen werden, werden
\Es{penetrierende Verletzungen} genannt. Stichverletzungen
durch Messer, Scheren oder ähnliche Gegenstände werden
zumeist mit einer geringeren Geschwindigkeit zugefügt als
Schusswaffenmunition und haben daher eine geringere sekundäre
Verletzungsfolge als Schussverletzungen. 

Stichwunden können oft sehr klein, schwach
blutend und daher harmlos aussehen. Um das Verletzungsausmaß richtig abschätzen
zu können, sollten folgende Aspekte beachtet werden:

\begin{itemize}
  \item Welche Körperteile sind betroffen?
  \item Wie viele Einstichwunden gibt es?
  \item Wie lang war die Klinge?
  \item Wie war der Einstichwinkel? (Welche Organe könnten noch betroffen sein?)
  \item Wurde nach dem Einstich das Messer (der Gegenstand) im Körper des
  Patienten umgedreht?
\end{itemize}

Im Körper steckende Messer, Stichwaffen
und andere \Es{Fremdkörper sollen niemals
präklinisch entfernt werden}! Eventuell verhindert ein noch steckender
Fremdkörper gerade noch eine unkontrollierte Blutung. Eine Entfernung darf
daher erst unter kontrollierten Bedingungen \Es{in einem Operationssaal}
erfolgen. Für den Transport muss man den \Es{Fremdkörper mittels eines
Verbandes sichern}, um ein Verrutschen und Folgeverletzungen zu verhindern.
}{
\begin{itemize}
  \item 
\end{itemize}
}
{./images/teaser/blut-tatort-hinterhof.jpg}
{}
{}



\begin{table}[H]
\Captionof{table}{Bilderserie:
\SlideHeading{Fremdkörper}}

\begin{tabularx}{\linewidth}[H]{XXX}
\addlinespace
\includegraphics[width=\linewidth]{./images/trauma/fremdkoerper-stich-fixierung.jpg}
\Captionof{figure}{\AuthNotesPicLegalNotOk{}{}Fixierung
eines penetrierenden Fremdkörpers.
\SourcePicture{Hauer}} 
&
\includegraphics[width=\linewidth]{./images/trauma/fremdkoerper-stich-messer-01.jpg}
\Captionof{figure}{\AuthNotesPicLegalNotOk{}{}Stich mit einem Küchenmesser. 
Das Messer verbleibt im Körper {\ldots}~
\SourcePicture{Gabriel}} 
&
\includegraphics[width=\linewidth]{./images/trauma/fremdkoerper-stich-messer-02.jpg}
\Captionof{figure}{\AuthNotesPicLegalNotOk{}{}{\ldots} bis in den Operationssaal.
\SourcePicture{Gabriel}} 
\\\addlinespace\bottomrule
\end{tabularx}
\end{table}

\subsection{Zusammenfassung}
Aus allen einsatztaktischen Überlegungen ergibt sich ein Bündel an Erstmaßnahmen,
welche Sie bei Unfällen durchführen müssen.

\MassnahmenNG{Spezielle Maßnahmen}{}{}{%
\begin{itemize}
  \item Sichern Sie die Unfallstelle ab!
  \item Verschaffen Sie sich einen Überblick (Wie viele Verletzte? ...) um ggf.
  weitere Einsatzkräfte nachzufordern!
  \item Erheben Sie den Verletzungsmechanismus! Stellen Sie (sich) folgende Fragen:
  \begin{itemize}
    \item Aus welcher Richtung hat die Kraft auf den Patienten eingewirkt?
    \item Wie schnell war das Fahrzeug unterwegs? bzw. Aus welcher Höhe ist der Patient gefallen?
    \item War der Patient angegurtet?
    \item Welche besondere Verformungen sind am Fahrzeug sichtbar?
    \item Wurde der Patient aus dem Fahrzeug geschleudert?
  \end{itemize}
\end{itemize}
}


\section{Wunden}\label{topic:wunden}\index{Wunden}

\Definition{Eine Wunde ist eine Gewebeschädigung der Haut
durch äußere mechanische, thermische oder chemische
Einwirkungen.}

\subsection{Einteilung von Wunden}

Wunden lassen sich \E{nach verschiedenen Gesichtspunkten} einteilen bzw.
unterscheiden.

Eine Einteilung bezieht sich auf die \Es{physikalische
Verletzungsursache}:

\begin{enumerate}
  \item \T{Mechanische Wunden}:
  \index{Wunden!mechanische}Entstehen durch spitze oder
  stumpfe Gewalteinwirkung auf den Körper (Stichwunde, Schürfwunde, Schusswunde, Bisswunde, ...)
  \item \T{Thermische Wunden}: \index{Wunden!thermische}
  Entstehen durch Kälte- oder Hitzeeinwirkung auf den Körper  (Verbrennungen,
  Stromeinwirkung, Erfrierung, ...)
  \item \T{Chemische Wunden}:
  \index{Wunden!chemische}Entstehen durch Verätzungen mit
  Laugen oder Säuren
\end{enumerate}
 
Eine weitere Unterteilung kann nach der \Es{Tiefe der
Schädigung} des Gewebes vorgenommen werden.

\begin{itemize}
  \item \Es{Oberflächliche Wunden:} Leichte Verletzungen,
  welche max. bis zur Lederhaut reichen.
  \item \Es{Tiefe Wunden:} Neben der Verletzung der Haut
  können hier auch Muskeln, Sehnen, größere Blutgefäße innere Organe und Knochen betroffen sein.
  \item \Es{Penetrierende Wunden:} tiefe Wunden mit
  Eröffnung mindestens einer Körperhöhle (Schädel, Brust, Bauch)
  % 	\item \Es{Geschlossene Wunden:} Verletzungen der
     % 	unteren Hautschichten ohne Eröffnung der Oberhaut. (Hämatom) % gestrichen
        % lt. STEUER
\end{itemize}





Unterschiedliche \E{Entstehungsmechanismen} erzeugen
unterschiedliche \E{Wundarten} bzw. \E{Wundformen}.

\Layout{56}{Wundarten}
{%

}{

}{
\begin{tabulary}{\linewidth}{LL}
\addlinespace
\noindent\Es{Schnittwunde}
&
Glatte Wundränder; tiefer liegende Nerven und Blutgefäße können verletzt sein
\\\addlinespace
\noindent\Es{Stichwunde}
&
Glatte Wundränder; meist kleine Wunde dafür aber tiefer
\\\addlinespace
\noindent\Es{Risswunde}
&
Unregelmäßige Wundränder; meist nur Schädigung der Haut und
	keine tieferliegenden Nerven und Blutgefäßen
\\\addlinespace
\noindent\Es{Quetschwunde}
&
Unregelmäßige Wundränder; es kann zum Aufplatzen der Haut kommen (Platzwunde)
\\\addlinespace
\noindent\Es{Rissquetschwunde} (RQW)
&
Kombination aus Riss- und Quetschwunde
\\\addlinespace
\noindent\Es{Schürf- und Kratzwunde} 
&
Oberflächliche Wunde
\\\addlinespace
\noindent\Es{Bisswunden} 
&
Kombination aus Riss-, Stich- und Quetschwunde
\\\addlinespace
\noindent\Es{Ablederung} 
&
Großflächiger Gewebedefekt, z.\,B. Skalpierung, große Lappenwunden
\\\addlinespace
\noindent\Es{Amputation} 
&
Abtrennung eines Körperteiles
\\\addlinespace
\noindent\Es{Pfählung~/ Einspießung} 
&
Eindringen und Verbleiben von spitzen aber z.\,T. unregelmäßigen Fremdkörpern
\\\addlinespace
\noindent\Es{Schusswunde} 
&
Meist klein aussehende Wunde, das eigentliche Schadensausmaß im Körperinneren
ist meist nicht ersichtlich, \Z{evtl. Ausschuss größer als
Einschuss}
\\\addlinespace
\noindent\Es{(Hämatom)} 
&
Bluterguss; Ansammlung von Blut im Gewebe außerhalb der Blutgefäße 
\\\addlinespace\bottomrule
\end{tabulary}
}{Wundarten.}{}


\begin{table}[H]
\begin{gWideParEnv}
\Captionof{table}{Bilderserie: \SlideHeading{Wunden}}

\begin{tabularx}{\linewidth}[H]{XXX}
\addlinespace
\includegraphics[width=\linewidth]{./images/wunde_rqw1.jpg}
\Captionof{figure}{\AuthNotesPicLegalOk{}{}Rissquetschwunde
vor der Versorgung im Spital.
\SourcePicture{Hauer}}
&
\includegraphics[width=\linewidth]{./images/wunde_rqw2.jpg}
\Captionof{figure}{\AuthNotesPicLegalOk{}{}Rissquetschwunde
nach der Versorgung im Spital.
\SourcePicture{Hauer}}
&
\includegraphics[width=\linewidth]{./images/pulsader_schnitt.jpg}
\Captionof{figure}{\AuthNotesPicLegalOk{}{}Schnittverletzung
oberhalb der "`Pulsader"'.
\SourcePicture{Hauer}}
\\\addlinespace
\includegraphics[width=\linewidth]{./images/trauma/stichwunde-klein.jpg}
\Captionof{figure}{\AuthNotesPicLegalNotOk{}{}Glatte
Wundränder und eigentlich ganz unauffällig: Die Stichwunde.
\SourcePicture{Hauer}} 
&
\includegraphics[width=\linewidth]{./images/pneumothorax_stichverletzung_diskret.jpg}
\Captionof{figure}{\AuthNotesPicLegalOk{}{}Stichverletzung:
unauffällig. \Ca{Dieser Patient ist
lebensgefährlich verletzt!}
(Pneumothorax)\SourcePicture{Hauer}} &
\includegraphics[width=\linewidth]{./images/trauma/stichwunde-thorax-tief-01.jpg}
\Captionof{figure}{\AuthNotesPicLegalNotOk{}{}Stichwunde
mit eröffneter Brusthöhle.
\SourcePicture{Hauer}} 
\\\addlinespace\bottomrule
\end{tabularx}
\end{gWideParEnv}
\end{table}




\begin{table}[H]
\begin{gWideParEnv}
\Captionof{table}{Bilderserie: \SlideHeading{Verätzungen und Verbrennungen}}

\begin{tabularx}{\linewidth}[H]{XXX}
\addlinespace
%\includegraphics[width=\linewidth]{./images/trauma/fremdkoerper-stich-fixierung.jpg}
% \Captionof{figure}{\AuthNotesPicLegalNotOk{}{}Fixierung
% eines penetrierenden Fremdkörpers.
% \SourcePicture{Hauer}} 
&
\includegraphics[width=\linewidth]{./images/trauma/burn-chem01.jpg}
\Captionof{figure}{\AuthNotesPicLegalNotOk{}{}Verätzung.
\SourcePicture{Hauer}} 
&
\includegraphics[width=\linewidth]{./images/trauma/burn-chem02.jpg}
\Captionof{figure}{\AuthNotesPicLegalNotOk{}{}Verätzung.
\SourcePicture{Hauer}} 
\\\addlinespace
\includegraphics[width=\linewidth]{./images/trauma/combustio-sa-01.jpg}
\Captionof{figure}{\AuthNotesPicLegalNotOk{}{}Großflächige
Verbrennung. \SourcePicture{Hauer}} 
&
\includegraphics[width=\linewidth]{./images/trauma/combustio-sa-02-hand.jpg}
\Captionof{figure}{\AuthNotesPicLegalNotOk{}{}Ablederung.
\SourcePicture{Hauer}} 
&
\includegraphics[width=\linewidth]{./images/trauma/combustio-sa-03-escharotomie.jpg}
\Captionof{figure}{\AuthNotesPicLegalNotOk{}{}Escharatomie.  Durch die
Ringförmige Verbrennung am Thorax zieht sich die Haut zusammen  und die Atmung
erschwert. Durch die "`Entlastungsschnitte"' kann der Patient wieder atmen. 
\SourcePicture{Hauer}} 
\\\addlinespace\bottomrule
\end{tabularx}
\end{gWideParEnv}
\end{table}


\begin{table}[H]
\begin{gWideParEnv}
\Captionof{table}{Bilderserie: \SlideHeading{Frakturen}}

\begin{tabularx}{\linewidth}[H]{XXX}
\addlinespace
\includegraphics[width=\linewidth]{./images/trauma/fraktur-offen-finger.jpg}
\Captionof{figure}{\AuthNotesPicLegalNotOk{}{}Nicht immer
spektakulär: Offene Fraktur des Mittelfingers.
\SourcePicture{Ben Stephenson, Cleveland, OH, USA.
Lizenz: CC-BY 2.0}} &
\includegraphics[width=\linewidth]{./images/trauma/fraktur-us-nativ.jpg}
\Captionof{figure}{\AuthNotesPicLegalNotOk{}{}Unterschenkelfraktur
nach Verkehrsunfall. 
\SourcePicture{Hauer}} 
&
\includegraphics[width=\linewidth]{./images/trauma/fraktur-us-roe.jpg}
\Captionof{figure}{\AuthNotesPicLegalNotOk{}{}Röntgenbild
der Unterschenkelfraktur. 
\SourcePicture{Hauer}} 
\\\addlinespace\bottomrule
\end{tabularx}
\end{gWideParEnv}
\end{table}
\subsection{Wundversorgung}

Die unterschiedlichen Wundarten benötigen auch unterschiedliche Wundversorgungen.
Bei den mechanischen Wunden wird die Blutstillung im Vordergrund stehen. Bei
thermischen Wunden wie z.\,B. Verbrennungen ist die Kühlung der betroffenen Stelle
als erstes durchzuführen. Bei chemischen Wunden sollte die Verdünnung der
ätzenden Substanz mit Wasser primär erfolgen. Schlussendlich werden aber alle
Wunden mit einer keimfreien Wundauflage zum Schutz vor Infektionen bedeckt.

Je nach Stärke der Blutung versorgen Sie Wunden mit einem Wundverband, einem
Druckverband oder in extremen Ausnahmesituationen durch eine Abbindung. Als
Erstmaßnahme bevor ein Verband zur Verfügung steht können Sie mit einer sterilen
Kompresse die Wunde abdrücken oder auch eine zuführende Arterie abdrücken.
Techniken \prettyref{topic:ehi-wunden}. 

\Layout{2}{Grundregeln der Wundversorgung}
{%
\label{topic:wundversorgung-grundregeln}%
\begin{itemize}
  \item Wunden dürfen nicht direkt berührt werden, außer es ist
  zur raschen Blutstillung unbedingt notwendig!
  \item Es muss immer mit Handschuhen und mit sterilem Material gearbeitet
  werden!~\footnote{Bei der Wundversorgung sind immer sterile Materialien zu
  verwenden. Papier-Handtücher o.\,ä. sind \E{nicht zulässig}! Eine Ausnahme
  kann lediglich bei der Versorgung des Stichs mittels eines sauberen
  Tupfers gemacht werden, da die Wunde minimal klein ist.}
  
  %\item Wunden sollen grundsätzliche \E{nicht} gereinigt
  %werden! Sie dürfen nur lose Glassplitter oder ähnliches
  %vorsichtig aus der Wunde entfernen.
  \item Wunden, die im Spital weiter versorgt werden, sollen
  nicht übermäßig gereinigt werden. \underline{Grobe}
  Verschmutzungen werden nur wenn unbedingt notwendig mit
  steriler Flüssigkeit \Z{(Kochsalzlösung (\ce{NaCl}-Lösung), Ringerlösung,
  o.\,ä.)} oberflächlich \underline{ab}gespült. Lose, kleine
  Fremdkörper wie z.\,B. Glassplitter dürfen entfernt werden.
  \item Nicht-weiterbehandlungspflichtige Wunden können
  mit einem nicht-alkoholischen
  Desinfektionsmittel\footnote{\Z{z.\,B.
  Octenisept\texttrademark{}}} gereinigt werden. Dabei erfolgt die Reinigung
  \Es{von Innen nach Außen} um keine weiteren Verunreinigungen in die Wunde
  zu bringen.
  \item Belassen Sie Fremdkörper (Messer, Schere oder sonstige
  Pfählungsgegenstände)  in der Wunde! Stabilisieren Sie diese Gegenstände mit 
  geeignetem Fixationsmaterial (z.\,B. Mullbinden).
  \item Ausgetretene Strukturen
  (z.\,B. Gehirn, Eingeweide) dürfen \E{nicht zurückgestopft} werden!
  Solche Organteile sollten mit steriler Kochsalzlösung (\ce{NaCl}-Lösung) feucht gehalten und anschließend
  steril  mit einem Verbandstuch abgedeckt werden!
  \item Ein steriler Wundverband mit der jeweiligen Situation
  angemessenem Material (z.\,B. Pflaster, Momentverband,
  Kompresse mit Mullbinde, Wundfolie, \ldots) soll angelegt werden (ausgenommen
  bei Bagatellverletzungen).
\end{itemize}
}{}{}{}{}

\begin{figure}[H]
\begin{center}
\includegraphics[width=0.5\linewidth]{./images/hirtler-aass/frische-Wunde-desinfizieren.pdf}

\Captionof{figure}{Desinfektion einer frischen Wunde. Wenn erforderlich, muss
\Es{von der Wunde weg} desinfiziert werden, um keine Keime in die Wunde zu
bringen. \SourcePicture{Hirtler.}}
\end{center}
\end{figure}


\subsection{Gefahren bei Wunden}
Jede Wunde hat (in unterschiedlichem Ausmaß) eine Auswirkung auf den
Gesamtorganismus. Je nach Wundart, stärke der Gewalteinwirkung, betroffener
Struktur und Sorgfalt bei der Wundversorgung gehen verschiedene Gefahren von
Wunden aus:
\begin{itemize}
  \item \Es{Verletzung wichtiger Gewebestrukturen (Organe):} Je nach
  betroffener Körperregion treten unterschiedliche Verletzungsmuster auf.  In
  den folgenden Abschnitten lernen Sie über spezielle traumatologische Notfälle
  wie z.\,B. Schädel-Hirn-Trauma, Thorax-Trauma, Bauchtrauma etc.
  \item \Es{Starke Blutung:} Bei großflächigen Wunden oder Verletzung wichtiger
  Blutgefäße kann es zu starken Blutungen und infolge dessen zu einem 
  Volumenmangelschock kommen. Blutstillung und Schockbekämpfung sind  die
  wichtigsten Maßnahmen.
  \item \Es{Starke Schmerzen:} Treten im Zusammenhang mit einer Verletzung 
  starke Schmerzen auf, muss ein Notarzt für eine Schmerztherapie gerufen werden.
  \item \Es{Eintritt von Krankheitserregern
  (Infektionsgefahr):} Jede Wunde ist eine
  Eintrittsstelle für Krankheitserreger in unseren Organismus. Der Patient
  sollte daher über einen aufrechten \Es{Tetanus}-Impfschutz verfügen 
  (s. \prettyref{topic:tetanus-impfschutz})\footnote{Zur Verifizierung eines
  aufrechten Tetanus-Impfschutzes muss jede Wunde einem Arzt vorgestellt
  werden.  Eine Impfung muss spätestens alle 10 Jahre aufgefrischt werden,  bei
  entsprechenden Verletzungen schon früher.\cite{BoehmerTaschRett6}}.  Bei
  Bisswunden kann über den Speichel einer erkrankten Tieres \Es{Tollwut}
  übertragen werden.  Versuchen Sie bei Bisswunden daher immer zu erfragen, wer der
  Besitzer des  Tieres ist und vermerken Sie dies auf Ihrem Einsatzprotokoll.
  Unter  Umständen kann es notwendig sein die Polizei einzuschalten.
\end{itemize}

%\subsection{Mechanische
%Wunden}\label{topic:wunden-meschanische}\index{Wunden!mechanische}





% \begin{tabularx}{\linewidth}[H]{XXX}
% \\\toprule
% \multicolumn{3}{c}{{\TabHeading{Diashow Muster}}}
% \\\midrule
% \includegraphics[width=\linewidth]{./images/stop.jpg}
% \Captionof{figure}{\AuthNotesPicLegalNotOk{}{}Muster
% \SourcePicture{}} 
% &
% \includegraphics[width=\linewidth]{./images/stop.jpg}
% \Captionof{figure}{\AuthNotesPicLegalNotOk{}{}Muster
% \SourcePicture{}} 
% &
% \includegraphics[width=\linewidth]{./images/stop.jpg}
% \Captionof{figure}{\AuthNotesPicLegalNotOk{}{}Muster
% \SourcePicture{}} 
% \\\midrule
% \includegraphics[width=\linewidth]{./images/stop.jpg}
% \Captionof{figure}{\AuthNotesPicLegalNotOk{}{}Muster
% \SourcePicture{}} 
% &
% \includegraphics[width=\linewidth]{./images/stop.jpg}
% \Captionof{figure}{\AuthNotesPicLegalNotOk{}{}Muster
% \SourcePicture{}} 
% &
% \includegraphics[width=\linewidth]{./images/stop.jpg}
% \Captionof{figure}{\AuthNotesPicLegalNotOk{}{}Muster
% \SourcePicture{}} 
% \\\bottomrule
% \end{tabularx}

%\subsubsection{Stichwunden}\label{topic:stichwunden}\index{Stichwunden}

% \begin{gWideParEnv}
% \begin{tabularx}{\linewidth}[H]{XXX}
% \\\toprule
% \multicolumn{3}{c}{{\TabHeading{Diashow}}}
% \\\midrule
%  &
% \includegraphics[width=\linewidth]{./images/stichverletzung_oberarm.jpg}
% \Captionof{figure}{\AuthNotesPicLegalOk{}{}Stichverletzung
% am Oberarm, Schwellung. \SourcePicture{Hauer}} &
% {leer}
% \\\bottomrule
% \end{tabularx}
% \end{gWideParEnv}


%\subsubsection{Schnittwunden}\label{topic:schnittwunden}\index{Schnittwunden}



%\subsubsection{Risswunden}\label{topic:risswunden}\index{Risswunden}

%\subsubsection{Quetschwunden}\label{topic:quetschwunden}\index{Quetschwunden}

%\subsubsection{Rissquetschwunden
%(RQW)}\label{topic:rqw}\index{Rissquetschwunden}\index{RQW}

% \begin{gDiasEnv3}
% \includegraphics[width=\linewidth]{./images/wunde_rqw1.jpg}
% \Captionof{figure}{\AuthNotesPicLegalOk{}{}Rissquetschwunde
% vor der Versorgung im Spital.
% \SourcePicture{Hauer}}
% &
% \includegraphics[width=\linewidth]{./images/wunde_rqw2.jpg}
% \Captionof{figure}{\AuthNotesPicLegalOk{}{}Rissquetschwunde
% nach der Versorgung im Spital.
% \SourcePicture{Hauer}}
% &
% \end{gDiasEnv3}





%\subsubsection{Schußwunden}\label{topic:schusswunden}\index{Schußwunden}




%\subsection{Thermische Wunden}\label{topic:wunden-thermische}\index{Wunden!thermische}

%\subsubsection{Brandwunden}\label{topic:brandwunden}\index{Brandwunden}

%\subsubsection{Erfrierungen}\label{topic:erfrierungen}\index{Erfrierungen}

%\subsection{Chemische Wunden}\label{topic:wunden-chemische}\index{Wunden!chemische}

%\subsubsection{Verätzungen}\label{topic:veraetzungen}\index{Verätzungen}


%%%%%%%%%%%%%%%%%%%%%%%%%%%%%%%%%%%%%%%%%%%%%%%%%%%%%%%%%%%
%%%%%%%%%%%%%%%%%%%%%%%%%%%%%%%%%%%%%%%%%%%%%%%%%%%%%%%%%%%

%\subsection{Stichverletzungen}

%NIEMALS  HERAUSZIEHEN!!!!!!!!

%steril abdecken + fixieren

%verletzte Strukturen??, dient ev. als {\quotedblbase}Stöpsel{\textquotedblleft}??

%Fremdkörper soll nicht tiefer rutschen



%%%%%%%%%%%%%%%%%%%%%%%%%%%%%%%%%%%%%%%%%%%%%%%%%%%%%%%%%%%
%%%%%%%%%%%%%%%%%%%%%%%%%%%%%%%%%%%%%%%%%%%%%%%%%%%%%%%%%%%
%%%%%%%%%%%%%%%%%%%%%%%%%%%%%%%%%%%%%%%%%%%%%%%%%%%%%%%%%%%
\clearpage
\section{Schädel-Hirn-Trauma (SHT)}
%Folge von Gewalteinwirkung auf den  Schädel
\index{Schädel-Hirn-Trauma|see{SHT}}

\Definition{Gewalteinwirkung auf den Kopf, die zusätzlich zu den nicht immer
vorhandenen Hautwunden und  zu den Frakturen des knöchernen Schädels
Funktionsstörungen und Verletzungen  des Gehirns hervorrufen.\cite{Gorgass7}}

(Vgl. Kap. Anatomie, \prettyref{topic:zns}.)





\Layout{60}{\TxBeschreibung{} und Arten}
{%
Die häufigsten Ursachen von Schädel-Hirn-Traumata (SHT) sind Verkehrsunfälle, 
Stürze und Gewaltverbrechen. Je nach Schwere und Mechanismus können
verschiedene  Formen des SHT auftreten. Man unterscheidet das 
\Es{gedeckte SHT}, bei dem die harte Hirnhaut (Dura mater) nicht eröffnet wird,
sowie das \Es{offene SHT}, bei dem die harte Hirnhaut eröffnet wird. Bei beiden
Arten kommt es zu Bewusstseinsstörungen. 

{\vspace{1cm}}

\Fazit{Das Leitsymptom des SHT
ist die \Es{Bewusstseinsstörung}.}
}{
\begin{itemize}
  \item Leitsymptom: \Es{Bewusstseinsstörung}.
  \item Arten:
  \begin{itemize}
    \item Gedecktes SHT: ohne  Eröffnung der harten Hirnhaut
    \item Offenes SHT: mit  Eröffnung der harten Hirnhaut
  \end{itemize}
\end{itemize}

}
{./images/teaser/imgp0722_00500.jpg}
{}
{GABS.}


\Layout{2}{Begleitende Verletzungen}
{%
\begin{itemize}
  \item Verletzung der Kopfschwarte  (Rissquetschwunde)
  \item Frakturen von
  \begin{itemize}
    \item Schädeldach
    \item Schädelbasis (Blutung bzw. Liquoraustritt
    aus Ohren oder Nase, Monokel-
    oder Brillenhämatom)
    \item Gesichtsschädel
    \item Unterkiefer
  \end{itemize}
\end{itemize}
}{

}{}{}{}




\Cave{Die Schwere eines SHT lässt sich nur bedingt von den äußerlich
sichtbaren  Verletzungen ableiten! Oft werden im Krankenhaus Frakturen
festgestellt,  ohne dass äußere Wunden vorliegen. Eine sorgfältige Untersuchung
ist daher bei jedem SHT Pflicht!}

\MassnahmenNG{Allgemeine Maßnahmen}{SHT}{\label{m:am-sht}}{

\begin{itemize}
  \item Vitale Bedrohung beurteilen. \par \TxBeiVit
  \TxMassVitBes
  \begin{itemize}
    \item \ce{O2} \VonBis{10}{15}\LMin
    \item Lagerung: s.\,u.
    \item Monitoring: Zusätzlich Beurteilung nach GCS (Verlauf beobachten und
    dokumentieren)\prettyref{topic:gcs}
  \end{itemize}
  \item Lagerung: leicht erhöhter Oberkörper, ca.
  30\,\textdegree{}, wenn keine Kontraindikation
  besteht.
  \item je nach Unfallmechanismus HWS-Immobilisation (z.\,B. mit Stifneck\texttrademark)
  \item Ggfs. Wundversorgung
  \begin{itemize}
    \item steriler Verband (steriles Abdecken)
    \item Knochen/Fremdkörper nicht entfernen/nicht zurück stopfen
  \end{itemize}
  %        \begin{itemize}
%            \item Gedecktes SHT:
%            \begin{itemize}
%                \item Steriler Verband
%            \end{itemize}
%            \begin{itemize}
%                \item steriler Verband locker auf offene
%                Verletzungen
%                \item Knochen / Hirnteile NICHT
%                zurückstopfen /  entfernen!
%                \item Ausgetretene Hirnmasse feucht steril abdecken
%            \end{itemize}
%        \end{itemize}
\end{itemize}
}

\subsection{Gedecktes SHT}

Ein gedecktes Schädel-Hirn-Trauma ist ein SHT ohne Eröffnung der harten
Hirnhaut. Je nach Schwere wird das SHT in drei Grade eingeteilt:

\begin{itemize}
\item \T{Gehirnerschütterung} -- SHT 1{\textdegree} --
Commotio cerebri
\item \T{Gehirnprellung} -- SHT 2{\textdegree} -- Contusio
cerebri
\item \T{Gehirnquetschung} -- SHT 3{\textdegree} 
\end{itemize}

\subsubsection{Gehirnerschütterung (Commotio cerebri) --
SHT 1\textdegree}
\index{Gehirnerschütterung}\index{Commotio (cerebri)}

\Layout{0}{\TxSymptome{} und Kennzeichen}
{%
Die Gehirnerschütterung ist die leichteste Form des
Schädelhirntraumas. Leitsymptom ist ein Trauma mit
anschließender kurzer Bewusstlosigkeit. Weiters kann es zu
flüchtigen neurologischen Ausfällen kommen, es kann auch eine
retrograde Amnesie (Gedächtnisverlust über den Vorgang
bzw. Unfall) bestehen. Eventuell klagt der Patient über
Übelkeit erbrechen und Schwindel. Die Gehirnerschütterung
zieht keine bleibenden Schäden nach sich.
}{
\begin{itemize}
    \item  Bewusstlosigkeit / Bewusstseinstrübung
    {\textless}\,1\,h
    \item Übelkeit, Erbrechen, Schwindel
    \item evtl. flüchtige neurologische Ausfälle
    \item retrograde  Amnesie  ({\quotedblbase}kann sich an den
    Unfallhergang nicht erinnern{\textquotedblleft} )
    \begin{itemize}
        \item im CT/MR 
        keine nachweisbaren Schäden
        \item keine  bleibenden Schäden
    \end{itemize}
\end{itemize}
}{}{}{}


%\MassnahmenNG{Spezielle Maßnahmen}{}{
%    \paragraph{Spezielle Maßnahmen beim SHT ---Gehirnerschütterung}
%    \begin{itemize}
%        \item 
%    \end{itemize}
%}


\subsubsection{Gehirnprellung (Contusio cerebri) --
SHT 2\textdegree}

\index{Gehirnprellung}\index{Contusio (cerebri)}

\Layout{0}{\TxSymptome{} und Kennzeichen}
{%
Die Gehirnprellung ist eine schwerere Form des Schädel-Hirntraumas. 
Leitsymptom ist die \Es{Bewusstlosigkeit größer
 als 1\,h}, beziehungsweise die Bewusstseinseintrübung
 >\,24\,h. Mittels weiterführender Diagnostik im Spital
 lassen sich die \Es{Hirnschädigungen} nachweisen. Es kann zu
 Nervenverletzungen kommen, dies kann zu vegetativen Störungen führen.
 Bleibende Schäden sind möglich.
}{
 \begin{itemize}
    \item Bewußtlosigkeit  {\textgreater}\,1\,h
    \item Bewußtseinstrübung {\textgreater}\,24\,h
    \item nachweisbare Gehirnschädigung / bleibende  Ausfälle
    \item ev. Gehirnnervenabrisse 
    \item  vegetative Störungen
\end{itemize}
}{}{}{}



%\MassnahmenNG{Spezielle Maßnahmen}{}{
%    \paragraph{Spezielle Maßnahmen beim SHT --- Gehirnprellung}
%    \begin{itemize}
%        \item 
%    \end{itemize}
%}




\subsubsection{Gehirnquetschung -- SHT 3{\textdegree}}
\index{Gehirnquetschung}

\Layout{0}{\TxSymptome{} und Kennzeichen}
{%
Die Gehirnquetschung ist eine sehr schwere Form des
Schädelhirntraumas. Es kommt zu einer \Es{Bewusstlosigkeit
länger als 24\,h} und zu teils sehr ausgeprägten
Begleitsymptomen wie zum Beispiel Hirndruckzeichen, Atem und
Kreislaufstörungen bis hin zum Atemstillstand durch Einklemmen
des Atemzentrums. 

Es kann weiteres zu intrakraniellen
Blutungen kommen (epidural, subdural, subarachnoidal).
Außerdem kommt es zu verschiedenen neurologischen Ausfällen
und Symptomen wie zum Beispiel zerebrale
Krampfanfälle. Mit \Es{bleibenden Schäden} und Langzeitfolgen ist
zu rechnen.
}{
\begin{itemize}
    \item Bewußtlosigkeit  {\textgreater}\,24\,h
    \item Atmungs-/ Kreislaufstörungen
    \item Hirndruck symptomatik
    \item evtl. Einklemmung des Atemzentrums ${\rightarrow}$
    zentraler Atemstillstand
    \item Behinderung der Hirndurchblutung
    \item Hirnblutungen
%     \begin{itemize}
%         \item epidural
%         \item subdural
%         \item subarachnoidal
%     \end{itemize}
    \item Lähmungen
    \item Zerebrale Krämpfe
\end{itemize}
}{}{}{}



%\MassnahmenNG{Spezielle Maßnahmen}{}{
%    \paragraph{Spezielle Maßnahmen beim SHT --- Gehirnquetschung}
%    \begin{itemize}
%        \item 
%    \end{itemize}
%}

\MassnahmenNG{Spezielle Maßnahmen}{SHT, gedeckt}{}{
%     \subsubsection{Spezielle Maßnahmen beim SHT --
%     }
    \begin{itemize}
        \item Allgemeine Maßnahmen beim SHT
        (\prettyref{m:am-sht})

    \end{itemize}
}


\subsection{Offenes SHT}

\Layout{2}{Hinweise}
{%
Hinweise auf ein offenes SHT sind:

\begin{itemize}
    \item große RQW mit tast-/sichtbaren  Knochensplittern
    \item Hirnaustritt (schlechte Prognose)
    \item Blutungen aus Nase/Ohren
    \item Liquoraustritt, oft mit Blut vermischt
    ({\quotedblbase}Tupfertest{\textquotedblleft})
\end{itemize}
}{

}{}{}{}


\MassnahmenNG{Spezielle Maßnahmen}{SHT, offen}{}{
%     \subsubsection{Spezielle Maßnahmen beim SHT -- }
    \begin{itemize}
        \item Allgemeine Maßnahmen beim SHT
        (\prettyref{m:am-sht})
        \item Ausgetretene Hirnmasse feucht steril abdecken
        und locker befestigen
    \end{itemize}
}


\clearpage
\subsection{Intrakranielle Blutungen}

\Layout{2}{\Ca{Luzides Intervall} (lichtes Intervall)}
{%
Die Symptome können eventuell erst Stunden nach der
Verletzung einsetzen. Dazwischen ist der Patient unauffällig.
Diese Zeitspanne wird \T{Luzides Intervall} genannt und ist
sehr gefährlich, da man den Patienten leicht falsch
einschätzen kann.
}{

}{}{}{}

\Layout{0}{Intrakranielle Blutungen}
{Intrakranielle Blutungen können sowohl in Folge eines SHT
als auch spontan (z.\,B. durch Platzen eines Aneurysmas)
auftreten. Aufgrund der inneren Blutung steigt der Druck im
Gehirn an, im Extremfall kann es zu einer
Hirnstammeinklemmung kommen. Patienten mit intrakranieller
Blutung sind vital bedroht! Beim SHT ist zu beachten, dass
die Symptome oft auch erst Stunden nach dem Unfall auftreten
können - siehe luzides Intervall. 
}{
\begin{itemize}
    \item stärkste Kopfschmerzen (langsam oder plötzlich einsetzend)
    \item Bewusstseinsstörungen bis Bewusstlosigkeit
    \item Pupillendifferenz, verlangsamte Pupillenreaktion
    \item evtl. Lähmungen oder Krämpfe
    \item evtl. Atemnot, unregelmäßige Atmung und Kreislaufstörung
    \item evtl. Druckpuls (Bradykardie mit hohen Blutdruckwerten als Folge des
    Hirndrucks)
    \item Übelkeit, Erbrechen
\end{itemize}
}{}{}{}



%\AuthNotes{Zusammenfassen, keine Unterteilung in Sub-,
%Epi- und Aubarachnoidalblutung, nur mehr "`Hirnblutung"'}


%\subsubsection{Epidurale Blutung} \index{Epidurale Blutung}

%\Definition{Blutung zwischen
%der harten Hirnhaut und dem Schädelknochen.}

%\Married{Bei der epiduralen Blutung kommt es zu einer Blutung zwischen
%der harten Hirnhaut und dem Schädelknochen, bedingt durch
%den Riss einer Hirnhaut-Arterie. Durch die Blutung steigt
%der Hrindruck, es kommt dabei zu \E{Hirndruckzeichen}.}
%{\begin{itemize}
%    \item Riß einer Meningealarterie
%    \item Einblutung zwischen Kalotte und Dura mater
%    \item {\textquotedblleft}lichtes{\textquotedblright} (luzides )
%    Intervall: einige Stunden
%\end{itemize}}



%\subsubsection{Subdurale  Blutung} \index{Subdurale Blutung}

%\Definition{Blutung
%unterhalb der harten Hirnhaut und oberhalt der weichen
%Hirnhaut.}

%\Married{Bei der subduralen Blutung kommend es zu einer ein Blutung
%unterhalb der harten Hirnhaut. Der Verlauf ist oft
%schleichend und kann sich über mehrere Tage erstrecken. Die
%genaue Unterscheidung zwischen epiduraler und subduraler
%Blutung ist präklinisch nicht möglich.}
%{\begin{itemize}
%    \item schleichender (!) Verlauf: oft über Tage
%    \item auch ohne  Trauma: Hypertoniker, Alkoholiker
%\end{itemize}}



%\subsubsection{Subarachnoidale Blutung -- SAB}
%\index{Subarachnoidale Blutung} (\index{SAB}SAB)

%\emph{Abkz.: \T{SAB}}



%\Definition{Bei der subarachnoidalen Blutung kommt es zu einer
%Blutung unterhalb der Spinnengewebeshaut (Arachnoidea).} 

%\Married{Dafür gibt es
%unterschiedliche Ursachen, sie kann sowohl durch ein Trauma
%ausgelöst werden als auch spontan entstehen. Bei der
%\E{spontanen} Subarachnoidalenblutung platzt typischerweise
%ein bereits bestehendes kleines Aneurysma eines Hirngefäßes und
%verursacht die Blutung. Dies kann aus heiterem Himmel
%passieren.

%\Fazit{Die subarachnoidalen Blutung kann durch ein Trauma
%oder spontan entstehen. Sie kann ein traumatologischer oder
%auch neurologischer Notfall sein.}

%Leitsymptom ist der plötzliche, schwere,
%\Es{"donnerschlagartige" Kopfschmerz}, welcher mit
%Hirndruckzeichen und einem oft raschen Verfall des
%Patienten einhergeht. Der Patient ist akut \E{vital
%bedroht!}}
%{
%\begin{itemize}
%    \item Trauma oder spontan
%    \item z.\,B. geplatztes Aneurysma, Brückenvenenabriß bei
%    Decelerationstrauma
%    \item Klassisch: plötzlicher, schwerer,
%    {\quotedblbase}donnerschlagartiger{\textquotedblleft} Kopfschmerz
%    \item evtl. Herd symptomatik, schneller
%    Verfall,
%\end{itemize}}





\subsection{Spezielle Diagnostik beim SHT}

Grundsätzlich ist beim Schädel-Hirn-Trauma die übliche
Diagnostik durchzuführen, die besonders beim SHT wichtigen
Punkte sind hier jedoch nochmal beschrieben.

Beurteilung von:
\begin{itemize}
	\item Vitalfunktionen
	\item Neurologie (Neurocheck)
	\item Verletzungen (Traumacheck)
	\item Unfallmechanismus (Anamnese)
\end{itemize}

\paragraph{Neurologische Beurteilung}

Die neurologische Beurteilung ist beim SHT natürlich
besonders wichtig und besonders genau durchzuführen.
Auffälligkeiten können Hinweise auf eine direkte
Hirnschädigung oder eine Hirnblutung sein. 

\Parallel{
\begin{description}
\item[Bewußtseinslage] klar/getrübt/bewusstlos
\item[Motorik] 
\item[Pupillenweite] Hirndruckzeichen
\item[Pupillenreaktion] Hirndruckzeichen
\end{description}

}{
\includegraphics[width=\linewidth]{./images/trauma/pupillendifferenz.jpg}
\Captionof{figure}{\Ca{Alarmzeichen:} Eine Pupillendifferenz kann ein wichtiger
Hinweis auf einen erhöhten Hirndruck sein!} }


\paragraph{Traumacheck} Ein Traumacheck ist obligat
und sollte keiner weiteren Erwähnung bedürfen (siehe
\prettyref{topic:traumacheck}).



\AuthNotes{GABS: Die einzelnen Massnahmen in Bezug auf die
unterschiedlichen SHTs gehören definiert.}

%\MassnahmenNG{Spezielle Maßnahmen}{}{
%    \paragraph{Spezielle Maßnahmen beim SHT --- Gehirnerschütterung}
%    \begin{itemize}
%        \item 
%    \end{itemize}
%}

\begin{itemize}
\item (Gefahr: hohe Querschnittslähmung)
\item Schädel/HWS/BWS bilden eine Linie
\item Jedes SHT ist im Spital abzuklären! (klinisch und  bildgebend)
\end{itemize}

\Cave{\Es{Jedes SHT ist im Spital abzuklären!} Es
besteht die Gefahr dass sich der Patient gegenwärtig in der
\Es{"`luziden Intervall"'} 
befindet und sich sein Zustand erst in einigen Stunden, dann aber schlagartig 
verschlechtert. Im Zweifelsfall wird der Patient dort zur  Beobachtung (24\,h
lang) aufgenommen.}



%%%%%%%%%%%%%%%%%%%%%%%%%%%%%%%%%%%%%%%%%%%%%%%%%%%%%%%%%%%
%%%%%%%%%%%%%%%%%%%%%%%%%%%%%%%%%%%%%%%%%%%%%%%%%%%%%%%%%%%
%%%%%%%%%%%%%%%%%%%%%%%%%%%%%%%%%%%%%%%%%%%%%%%%%%%%%%%%%%%

\clearpage
\section{Wirbelsäulentrauma}
\label{topic:wirbelsaeulentrauma}

\AuthNotes{KOCH: Überarbeiten, evt etwas kürzen.}

\Definition{Gewalteinwirkung auf die Wirbelsäule, die zur Verschiebung oder
zur  Fraktur von Wirbeln mit oder ohne Rückenmarkschädigung führt.\cite{Gorgass7}}

\paragraph{Typische Unfallmechanismen}

\begin{itemize}
  \item Sturz aus großer Höhe (> 3 m)
  \item Motorradunfall
  \item eingeklemmte Personen bei VU
  \item Schlag auf Kopf (Stauchung der HWS)
  \item Schleuderbewegung bei VU, wenn Patient nicht angeschnallt war  oder
  Kopfstütze fehlt
  \item Suizidversuch durch Erhängen
\end{itemize}

\paragraph{\TxGefahren}

\begin{itemize}
  \item Kompression des Rückenmarks (aufgrund von Einblutungen ins  Rückenmark 
  und dadurch bedingtes Ödem)
  \item Durchtrennung des Rückenmarks
  \item Lähmungen (Parese)
  \item Sensibilitätsstörungen (Parästhesie)
\end{itemize}






\MassnahmenNG{Spezielle Maßnahmen}{Wirbelsäulentrauma --- Allgemein}{}
{%
\begin{itemize}
  \item bei ansprechbaren Patienten: immer HWS-Immobilisation (Stifneck)
  \item Rettung mit Schaufeltrage oder KED-System
  \item Lagerung auf/Schienung mit Vakuummatratze (wenn Schocklagerung
  notwendig,  Tragentisch schräg stellen)
  \item ev. Sauerstoffgabe
  \item bei Querschnittslähmung: Kennzeichnung der Sensibilitätsgrenze mit Angabe der Uhrzeit
  \item ev. Wundversorgung
  \item ev. Schockbekämpfung
  \item regelmäßige Kontrolle der Atmung
  \item Erheben und Dokumentieren des Unfallmechanismus
\end{itemize}
}

\subsection{Schleudertrauma, Peitschenschlag}

Schleudertraumen kommen bei Auffahrunfällen vor, insbesondere dann, wenn der
Patient  nicht angeschnallt war, oder die Kopfstütze fehlt. Es kommt zu einer
Zerrung  des Halsmarks, welche sich durch
\begin{itemize}
  \item Schmerzen im HWS-Bereich
  \item Kopfschmerzen
  \item Schwindel, Tinnitus
  \item ev. neurologische Ausfälle (in schweren Fällen)
\end{itemize}
bemerkbar macht.

\subsection{Querschnittssyndrom}

\Layout{0}{}{%
Kommt es zu einer Druckerhöhung (z.\,B. aufgrund
einer Blutung) im Rückenmark oder zu einer Durchtrennung des
Rückenmarks können Symptome eines Querschnitts auftreten. Ein
sorgsamer und vorsichtiger Umgang bei der Rettung, der
Untersuchung und dem Transport des Patienten ist unbedingt
notwendig.
}{
\begin{itemize}
  \item Ausfall von Motorik und/oder Sensibilität
  \item Blutdruckabfall aufgrund fehlender Gefäßregulation
  \item Schlaffe oder spastische Lähmung der Extremitäten
  \item Inkontinenz aufgrund von Blasen-/Darmlähmung
  \item Dauererektion (Priapismus)
\end{itemize}
}{}{}{}

\subsection{Spinaler Schock}

\Layout{0}{\TxBeschreibung}
{Aufgrund des Unfalls kommt es zu einer plötzlichen
Durchtrennung des Rückenmarks, wodurch die Gefäßregulation (und somit der Gefäßtonus) versagt.
Somit sinkt der Blutdruck rasch ab, es kommt zu einer lebensbedrohlichen
Kreislaufstörung.
}{
\begin{itemize}
	\item deutlicher Blutdruckabfall aufgrund fehlender Gefäßregulation
	\item Querschnittslähmung	
	\item Atemlähmungen (bei hohem Querschnitt (im Bereich der HWS))
\end{itemize}
}{}{}{}

%%%%%%%%%%%%%%%%%%%%%%%%%%%%%%%%%%%%%%%%%%%%%%%%%%%%%%%%%%%
%%%%%%%%%%%%%%%%%%%%%%%%%%%%%%%%%%%%%%%%%%%%%%%%%%%%%%%%%%%
%%%%%%%%%%%%%%%%%%%%%%%%%%%%%%%%%%%%%%%%%%%%%%%%%%%%%%%%%%%

\section{Thoraxtrauma}

\Definition{Verletzung des Brustkorbes und/oder der darin
enthaltenen Organe.}

Thoraxverletzungen müssen als besonders gefährlich eingestuft
werden, da im Brustkorb (Thorax) lebenswichtige Organe (Herz,
Lunge) und große Gefäße liegen.  Nahezu jeder 2. Verkehrstote
ist Opfer eines Thoraxtraumas. \cite{Gorgass7}

\Cave{Sie sollten niemals aufgrund der äußeren Verletzungen auf die  Schwere
des Thoraxtraumas schließen, da schwere innere Verletzungen auftreten können!}

\paragraph{\TxSymptome des Thoraxtraumas}~

\begin{tabularx}{\linewidth}{XX}
\TabHeading{Allgemeine Symptome}
&
\TabHeading{Spezifische Symptome}
\\\midrule
\begin{itemize}
  \item Dyspnoe
  \begin{itemize}
    \item Tachypnoe (beschleunigte Atmung)
    \item Zyanose
  \end{itemize}
  \item atemabhängige  Schmerzen,  Inspirationsstop
  \item ev. Hautemphysem (Luftansammlung in der Haut, typisches
  {\quotedblbase}Knistern{\textquotedblleft})
  \item Prellmarken (z.\,B. Lenksäule von PKW)
  \begin{itemize}
    \item gründlicher Traumacheck!
  \end{itemize}
\end{itemize}
&
\begin{itemize}
  \item instabiler Thorax
  \begin{itemize}
    \item Sternumfraktur
    \item Serienrippenfraktur (zus. paradoxe Atmung --  Brustkorb senkt sich beim
    Einatmen)
  \end{itemize}
  \item Dyspnoe + Prellmarken +  Bluthusten
  \begin{itemize}
    \item Lungenprellung: Gasaustausch, da  Lungenabschnitte
    geschädigt
  \end{itemize}
  \item Brustschmerz, Schock
  \begin{itemize}
    \item Rhythmusstörungen ${\rightarrow}$ Verletzung des  Herzens
    \item schwerer Schock, schneller Verfall ${\rightarrow}$ Riss der  Aorta
    \item Blut im Schwall aus dem Mund ${\rightarrow}$
    Speiseröhrenverletzung
  \end{itemize}
\end{itemize}
\\
\end{tabularx}



\subsection{Pneumothorax}\label{topic:pneumothorax}\index{Pneumothorax}

\Definition{Luftansammlung im Pleuraraum}

(Vgl. Kap. Anatomie, \prettyref{topic:atemmechanik} und
\prettyref{topic:pleura})
\index{Pneumothorax!geschlossener}
\index{Pneumothorax!offener}
\index{Pneumothorax!Spontan-}
\index{Pneumothorax!Spannungs-}
\index{Spontanpneumothorax}
\index{Spannungspneumothorax}

Zwischen Lungenoberfläche und innerer Thoraxwand herrscht
beim Gesunden ein Unterdruck gegenüber dem äußeren Luftdruck.
Die Lunge ist dadurch flexibel im Brustkorb aufgespannt. Beim
Pneumothorax gelangt durch eine Verletzung Luft in diesen
Pleuraspalt, welche den Unterdruck aufhebt und die Lunge
zieht sich zusammen. 

Man unterscheidet zwischen \T{offenem
Pneumothorax} (Luft strömt durch eine Verletzung der
Brustwand von außen in den Pleuraspalt) und
\T{geschlossenem Pneumothorax} (Luft strömt durch eine
Verletzung der Lungenoberfläche von innen in den
Pleuraspalt). 

Entsteht durch die Form der Wunde ein
Ventilmechanismus wird beim Einatmen Luft in den Pleuraspalt
gesogen, während sich beim Ausatmen die Wunde verschließt und
die Luft nicht mehr entweichen kann. Dadurch füllt sich der
Pleuraspalt mehr und mehr mit Luft, bis schließlich im
betroffenen Pleuraspalt ein Überdruck entsteht. Im Extremfall
drückt der Überdruck den Mittelfellraum auf die Gegenseite,
wodurch auch noch die gesunde Lungenhälfte und ev. die
Hohlvene komprimiert wird. In diesem Fall spricht man vom
\T{Spannungspneumothorax}.

Ein Pneumothorax kann auch ohne vorangegangene Verletzung auftreten. Diese Form
des Pneumothorax nennt man \T{Spontanpneumothorax} und tritt am häufigsten bei
jungen Menschen auf. Symptome des Spontanpneumothorax sind Atemnot und
Todesangst.



%\begin{itemize}
%   \item Lungen- (innerer) oder  Brustwandverletzung (äußerer)
%    \begin{itemize}
%        \item Luft gelangt in den Pleuraspalt
%        \item Unterdruck aufgehoben
%        \item Lunge kollabiert auf betroffener Seite
%    \end{itemize}
%    \item "`unkomplizierter"' Pneumothorax
%    \begin{itemize}
%        \item Luft kann ungehindert in Pleuraspalt und  wieder heraus
%        \item {\quotedblbase}nur{\textquotedblleft} Ausfall des betroffenen
%        Lungenflügels
%    \end{itemize}
%    \item Ventil-, \index{Spannungspneumothorax}Spannungspneumothorax
%    \begin{itemize}
%        \item Luft kann zwar hinein, aber nicht wieder  heraus
%        \item Druck auf betroffener Seite steigt ${\rightarrow}$  Organe (Lunge,
%        Herz) werden auf die andere Seite gedrückt
%        \begin{itemize}
%            \item gesunde Seite wird verdrängt + Abknickung  großer
%            Gefäße!
%        \end{itemize}
%    \end{itemize}
%\end{itemize}


%Therapie

%Pneumothorax, unkompliziert:

\begin{table}[h]
\begin{gWideParEnv}
\Captionof{table}{Bilderserie \SlideHeading{Pneumothorax}.}

\begin{tabularx}{\linewidth}[H]{XXX}
\includegraphics[width=\linewidth]{./images/pneumothorax_xray_halb.jpg}
\Captionof{figure}{\AuthNotesPicLegalNotOk{}{}Pneumothorax. 
Beachte die Randlinie zwischen Lunge
(unten, innen) und dem luftgefüllten
Pleuraspalt (oben, außen) \SourcePicture{Hauer}} &
\includegraphics[width=\linewidth]{./images/pneumothorax_xray_endstadium.jpg}
\Captionof{figure}{\AuthNotesPicLegalNotOk{}{}Kompletter
Pneumothorax der rechten Seite. Die Lunge ist kollabiert.  
\SourcePicture{Hauer}} 
&
\includegraphics[width=\linewidth]{./images/pneumothorax_stichverletzung_diskret.jpg}
\Captionof{figure}{\AuthNotesPicLegalOk{}{}Diese
unauffällige Stichverletzung hat einen Pneumothorax
verursacht.
\Ca{Dieser Patient ist lebensgefährlich verletzt!} 
\SourcePicture{Hauer}} \\\bottomrule
\end{tabularx}
\end{gWideParEnv}
\end{table}

\MassnahmenNG{Spezielle Maßnahmen}{Pneumothorax}{}
{
\begin{itemize}
  \item \TxMassVitBes
  \begin{itemize}
    \item \ce{O2}: 10\,--\,15\LMin
    \item Lagerung: Oberkörper hoch. Wenn stabile
    Seitenlage notwendig ist dann Lagerung auf die
    \E{verletzte} Seite, damit die unverletzte
    Lunge  nicht unnötig belastet wird.
  \end{itemize}
  \item Bei offener Thoraxverletzung: \Es{Nicht
  luftdicht verbinden!} Es kann sonst ein
  Spannungspneumothorax entstehen. Steril abdecken.
\end{itemize}
}

\Z{Der Notarzt kann mittels Drainage einen Pneumothorax
entlasten.}

%\AuthNotes{KOCH: 

%Es fehlt:
%Serienrippenfraktur erklären
%Paradoxe Atmung
%Gefahr von Blutverlust
%Therapie
%Ev. auch kurz etwas über STernumfraktur schreiben?}

\subsection{Rippenfraktur, Serienrippenfraktur und instabiler Thorax}
\label{topic:serienrippenfraktur}
\index{Serienrippenfraktur}

\Layout{0}{\TxBeschreibung}
{%
Von einer \Es{Serienrippenfraktur} spricht man,
wenn mindestens drei benachbarte Rippen gebrochen sind. Sind
die Rippen an verschiedenen Stellen gebrochen oder gar
zertrümmert, werden diese Teile bei der Einatmung nach innen
gezogen und bei der Ausatmung nach außen gedrückt. Man
spricht in diesem Fall von \Es{paradoxer Atmung} bzw.
\Es{instabilem Thorax}, welche eine starke Beeinträchtigung
der Lungenfunktion nach sich zieht. Die Folge ist hochgradige
Atemnot. Außerdem treten (z.\,T. heftige) Schmerzen auf, die
sich bei der Einatmung verstärken. Rippenbruchstücke können
darunter liegende Organe verletzen. Je nach Rippe(n) kann es
zu Lungeneinrissen, Leber- und Milzverletzungen mit starken
inneren Blutungen kommen. Der Blutverlust kann \VonBis{2}{3} Liter
betragen \cite{LPN3}. In diesem Fall herrscht höchste
Schockgefahr.
}{
\begin{itemize}
  \item (z.T. schwere) Atemnot, Zyanose, Tachypnoe
  \item (z.T. heftige) Schmerzen beim Einatmen
  \item ev. Schocksymptome
  \item ev. paradoxe Atmung
\end{itemize}
}{}{}{}

\subsection{Sternumfraktur}\label{topic:sternumfraktur}\index{Sternumfraktur}
Eine Fraktur des Brustbeins (Sternum) tritt vor allem bei nicht angeschnallten
Fahrzeuglenkern auf, wenn der Brustkorb gegen das Lenkrad prallt. Meistens treten
starke Schmerzen auf, gelegentlich lassen sich eine Delle oder eine Stufe tasten.

\subsection{Verletzung des Herzens und der großen Gefäße}
Verletzungen des Herzens und der großen Gefäße treten am wahrscheinlichsten bei
Stürzen aus großer Höhe und Verkehrsunfällen mit Frontalaufprall und großer
Geschwindigkeit auf, seltener bei Schuss- oder Stichverletzungen. Folgende
Verletzungen sind typisch:

\begin{itemize}
  \item \Es{Aortenruptur}: Durch einen Einriss der Aorta kommt es zu schweren inneren Blutungen.
  \item \Es{Herzkontusion}: Hier kommt es durch Quetschungen der Herzmuskulatur
  zu Sickerblutungen. Bei Sternumfrakturen mit Herzrhythmusstörungen sollte man
  an eine Herzkontusion denken.
  \item \Es{Herzbeuteltamponade}: Es kommt zu einer Einblutung in den
  Herzbeutel,  wodurch das Herz komprimiert wird und sich nicht mehr mit Blut
  füllen  kann. Leitsymptome sind gestaute Halsvenen und schwerer Schockzustand
  (kardiogener Schock).
\end{itemize}

\section{Bauch- und Beckentrauma}
\index{Bauchtrauma}

Ein Bauchtrauma oder Abdominaltrauma ist eine Verletzung des Abdomens. In einigen
Büchern ist das Bauchtrauma als Untergruppe des akuten Abdomens klassifiziert.
Häufigste Ursachen sind Verkehrsunfälle, Stürze aus großer Höhe sowie Stoß- und
Schlagverletzungen. \cite{RedelsteinerHRNS1} Grundsätzlich ist zwischen offenem
(perforierendem, penetrierendem) und geschlossenem (stumpfem) Bauchtrauma zu
unterscheiden.

\begin{table}[H]
\Captionof{table}{Offenes und geschlossenes Bauchtrauma.}

\begin{tabularx}{\linewidth}{XX}
\TabHeading{Offenes (perforierendes) Bauchtrauma}
&
\TabHeading{Geschlossenes (stumpfes) Bauchtrauma}
\\\midrule
\begin{itemize}
  \item  Eröffnung der Bauchhöhle, ev. Austritt von Darmschlingen
  \item Stich-, Pfählungsverletzung
  \item Schnittverletzung, {\quotedblbase}Platzbauch{\textquotedblleft}
\end{itemize}
&
\begin{itemize}
  \item Schlag, Tritt
  \item Prellung durch Lenkstange bei PKW-Unfall
\end{itemize}
\\\bottomrule
\end{tabularx}
\end{table}


\subsection{Offenes Bauchtrauma}

\Layout{20}{\TxBeschreibung}{
Offene Bauchtraumata sind seltener als geschlossene
Bauchtraumata. Beim offenen Bauchtrauma kann es zum Austritt
von Darmschlingen kommen. Die Größe der Eintrittswunde sagt
nichts über den Schweregrad der Verletzung aus, da unter der
Eintrittswunde unterschiedlich schwere Organschäden mit
inneren Blutungen versteckt sein können.
}{   
\TODO{GABS}{2010-04-29}{Punkte fehlen}{Punkte fehlen}
}
{./images/trauma/bauchtrauma-offenes-01.jpg}
{Offenes Bauchtrauma mit Austritt von Darmschlingen.}
{Hauer}


\MassnahmenNG{Spezielle Maßnahmen}{offenes Bauchtrauma}{}{

\begin{itemize}
  \item \TxMassVitBes
  \begin{itemize}
    \item Lagerung: Bauchdecke entspannen , Knierolle
  \end{itemize}
  \item Darmteile nicht  zurückstopfen ${\rightarrow}$ sonst
  Darmriß/Darmverschluß
  \item steril abdecken, mit Ringerlösung feucht halten
  \begin{itemize}
    \item Gefahr des Eindringens von Keimen  (Peritonitis/Sepsis!)
  \end{itemize}
\end{itemize}
}



\subsection{Geschlossenes Bauchtrauma}

Geschlossene Bauchtraumata sind aufgrund der inneren Blutungen besonders
gefährlich. So kann es z.\,B. bei einer Milzruptur zu einem Blutverlust von 4 Liter kommen!

\Fazit{Immer nach \Es{Prellmarken} suchen. \Es{Patienten
entkleiden!}}

\Layout{2}{\TxSymptome}
{%
\begin{itemize}
  \item starke  Bauchschmerzen
  \item brettharte  Bauchdecke  (Abwehrspannung, D\'efense )
  \item Leitsymptom: Schockzeichen
  \item beim Bewußtlosen schwierig zu  diagnostizieren
  \begin{itemize}
    \item gründlich nach \index{Prellmarken}\Es{Prellmarken}
    suchen!
  \end{itemize}
\end{itemize}
}{
\begin{itemize}
  \item starke  Bauchschmerzen
  \item brettharte  Bauchdecke  (Abwehrspannung, D\'efense )
  \item Leitsymptom: Schockzeichen
  \item beim Bewusstlosen schwierig zu  diagnostizieren
  \begin{itemize}
    \item gründlich nach \index{Prellmarken}\Es{Prellmarken}
    suchen!
  \end{itemize}
\end{itemize}
}{}{}{}


\Layout{2}{\TxGefahren}{%
\begin{itemize}
  \item Gefäßverletzungen
  \begin{itemize}
    \item Aorta, Vena cava, \ldots{}
    \item massive innere Blutungen ${\rightarrow}$ Volumenmangelschock
  \end{itemize}
  \item Zerreißung  von Organen (\Es{Milz},
  \Es{Leber})
  \begin{itemize}
    \item massive \Es{innere Blutungen} (bis über 4 Liter) ${\rightarrow}$
    Volumenmangelschock
  \end{itemize}
  \item Zerreißung  von Hohlorganen (Magen, Darm,Gallengänge )
  \begin{itemize}
    \item Austritt von Speisebrei, Magensaft, Stuhl, Gallensekret
    \begin{itemize}
      \item Bauchfellentzündung (Peritonitis) nach Std. bis  Tagen
      \item Sepsis (hohe Letalität), septischer Schock
    \end{itemize}
  \end{itemize}
\end{itemize}
}{
\begin{itemize}
  \item Gefäßverletzungen
  \begin{itemize}
    \item Aorta, Vena cava,...
    \item massive innere Blutungen ${\rightarrow}$ Volumenmangelschock
  \end{itemize}
  \item Zerreißung  von Organen (\Es{Milz},
  \Es{Leber})
  \begin{itemize}
    \item massive \Es{innere Blutungen} (bis über 4 Liter) ${\rightarrow}$
    Volumenmangelschock
  \end{itemize}
  \item Zerreißung  von Hohlorganen (Magen, Darm,Gallengänge )
  \begin{itemize}
    \item Austritt von Speisebrei, Magensaft, Stuhl, Gallensekret
    \begin{itemize}
      \item Bauchfellentzündung (Peritonitis) nach Std. bis  Tagen
      \item Sepsis (hohe Letalität), septischer Schock
    \end{itemize}
  \end{itemize}
\end{itemize}
}{}{}{}








\subsubsection{Milzruptur}\index{Milzruptur}

\Layout{0}{}{%
Die Milz ist ein "`teuflisches"' Organ. Sie besitzt rund um
das gut durchblutete Gewebe eine Bindegewebskapsel. Das
Schlimme daran: Das Gewebe kann reißen, die Kapsel bleibt
jedoch intakt. Es kommt zuerst nur zu einer geringfügigen
Einblutung, dem Patienten geht es zunächst gut. Nach
einiger Zeit reißt dann die Kapsel und es kommt zu einer
ungehemmten Blutung in den Bauchraum. \Es{Der Patient
verfällt von einem Moment auf den anderen in einen Volumenmangelschock.}
}{
\begin{itemize}
  \item Schmerzen im linken Oberbauch
  \item manchmal Ausstrahlung der Schmerzen in die linke Schulter (Kehr-Zeichen)
  \item Volumenmangelschock
\end{itemize}
}{}{}{}

\paragraph{Primäre (einzeitige) Milzruptur} Die Milz und die Kapsel
reißen gleichzeitig, es kommt zu einem unmittelbar merkbaren
Schockzustand.

\paragraph{Sekundäre (zweizeitige) Milzruptur} Die Kapsel reißt erst
einige Zeit (Stunden bis Tage) nach dem Riss des Milzgewebes
\E{(Latenzzeit)}. Es kommt zu einem \Es{zeitlich versetztem} und oft
\Es{unerwartetem Schockgeschehen}. Mitunter reißt die
Kapsel erst Stunden nach dem eigentlichen Unfall. \Es{Daher
ist es wichtig jeden Patienten, bei dem auch nur der geringste
begründete Verdacht auf eine Milzverletzung besteht, einer
weiterführenden Untersuchung im Spital zuzuführen!}
\Z{Mittels Ultraschall kann auf eine Milzblutung untersucht
werden.}

\Cave{Cave! {\quotedblbase}\index{Zweizeitige
Milzruptur}Zweizeitige Milzruptur{\textquotedblleft}}

\paragraph{Zusammenfassung:}
\begin{itemize}
  \item primäre Milzruptur: sofort, Kapsel durchgerissen
  \item sekundäre Milzruptur:
  \begin{itemize}
    \item Bei Unfall zunächst nur Riss des Milzgewebes, Kapsel noch intakt.
    \item symptomfreies Intervall (Latenzzeit), Sickerblutung
    \item Später reißt auch die Kapsel  ${\rightarrow}$ Schock!
  \end{itemize}
  %Sickerblutung, Kapsel anfangs noch intakt
  ${\rightarrow}$ Latenz bis Riss der Kapsel
\end{itemize}




\subsubsection{Weitere innere Verletzungen}

Der Bauchraum beherbergt viele, teilweise sehr gut
durchblutete Organe die bei einer Verletzung zu erheblichen
Blutverlusten führen können. Präklinisch äußern sich diese
Verletzungen im durch den Blutverlust hervorgerufenen
Schockzustand. Weitere Hinweiszeichen sind Prellmarken oder
Schmerzen in der entsprechenden Gegend.

Grundsätzlich kann man ursächlich gegen diese Verletzung
wenig tun, im Vordergrund steht die Erkennung,
Schockbehandlung und der zügige Transport in ein
Traumazentrum.


\Layout{0}{Leberruptur}
{%
Die Leber ist auch gut durchblutet und kann im Falle einer Verletzung große
Probleme machen. Ursachen sind meist stumpfe Verletzungen, z.\,B. ein
Lenkradaufprall oder eine Sicherheitsgurtkompression
}{
\begin{itemize}
  \item Druckschmerz im rechten Oberbauch
  \item Volumenmangelschock
  \item späterer Kapselriss möglich (zweizeitige Ruptur wie bei Milz, siehe oben)
\end{itemize}
}

\paragraph{Aortenruptur}\index{Aortenruptur}\label{topic:aortenruptur}Durch
einwirkende starke Kräfte kann es zu einem Riss in der Aorta
kommen. Besonders bei Stürzen aus großen Höhen kann dies
vorkommen. Da es sich bei der Aorta um das größte Gefäß des
Körpers handelt ist die Prognose äußerst schlecht.

\paragraph{Nierenruptur}\index{Nierenruptur}\label{topic:nierenruptur}
Die Niere als naturgemäß gut durchblutetes Organ kann im
Falle einer Verletzung auch massive Probleme machen.


\subsection{Beckentrauma}
\label{topic:beckentrauma}

Ein Beckentrauma ist eine Verletzung des Beckens. In manchen Lehrbüchern werden
sie zu den Bauchverletzungen gezählt. 

Besonders gefährlich sind Frakturen des Beckengürtels
(\prettyref{topic:beckenguertel}). Stehen und Gehen ist dann nicht mehr möglich,
der Patient verspürt starke Schmerzen bei jeder passiven Bewegung \E{beider}
Beine. Da der Beckenknochen gut durchblutet ist, können binnen weniger Minuten
bis zu fünf Liter Blut (also fast die gesamte Blutmenge) verloren gehen. Es besteht höchste
\Es{Schockgefahr!} Eine \Es{Kompression des Beckens} ist \Es{überlebenswichtig}
um den möglichen großen Blutverlust einzudämmen.

\begin{itemize}
  \item Entstehung durch erhebliche Gewalteinwirkung  (Verschüttung, Überrolltrauma)
  \item massive  Blutungen (bis 5\,l in wenigen  Minuten),
  Gefahr des Ausblutens (ev. ohne sichtbare  Blutungsquelle!)
  \item Mitverletzung von Organen im kleinen  Becken:
  \begin{itemize}
    \item Dickdarm ${\rightarrow}$ Blutung aus Anus
    \item Harnsystem ${\rightarrow}$ Blut im Harn, Harn in der
    Bauchhöhle
    %\item (Ausnahme: Spontanruptur der gefüllten
    %Harnblase bei stumpfem Bauchtrauma)
    \item Hoden ${\rightarrow}$ Hämatome am Skrotum und am  Damm
  \end{itemize}
\end{itemize}


\MassnahmenNG{Spezielle Maßnahmen}{instabiles
    Beckentrauma}{}
{
\begin{itemize}
  \item \TxMassVitBes
  \begin{itemize}
    \item Lagerung: flach
    \item \ce{O2}: 10-15\,l
  \end{itemize}
  \item Schockbekämpfung
  \item Wirbelsäulenschonende, situationsgerechte Rettung (z.\,B. mittels
  Schaufeltrage)
  \item Lagerung auf/Schienung mit Vakuummatratze (wenn Schocklagerung notwendig, Tragentisch schräg stellen)
  \item Bei Beckenfraktur:
  \index{Beckengurt}\Es{Beckengurt}, zur Not mit
  Gürtel
\end{itemize}
}


%%%%%%%%%%%%%%%%%%%%%%%%%%%%%%%%%%%%%%%%%%%%%%%%%%%%%%%%%%%
%%%%%%%%%%%%%%%%%%%%%%%%%%%%%%%%%%%%%%%%%%%%%%%%%%%%%%%%%%%
%%%%%%%%%%%%%%%%%%%%%%%%%%%%%%%%%%%%%%%%%%%%%%%%%%%%%%%%%%%
\clearpage
\section{Polytrauma}
\label{topic:polytrauma}

%\AuthNotes{Neufassung}

\Definition{Gleichzeitige \Es{Verletzung mehrerer Körperregionen} oder
Organsystemen, von denen \Es{wenigstens eine} allein
\Es{oder die Kombination mehrerer} \Es{lebensgefährlich}
ist.}

Beispiele:
\begin{itemize}
  \item Bauchtrauma + Oberschenkel; SHT + WS + Thoraxtrauma etc.
  \item gebrochener Finger + gebrochene Zehe ${\rightarrow}$ kein
  Polytrauma
\end{itemize}

Häufigkeit verletzter Körperregionen bei Polytraumatisierten (nach
\cite{BoehmerTaschRett6} und \cite{Gorgass7}):
\begin{itemize}
  \item 60-67\,\% Schädel-Hirn
  \item 45-75\,\% Extremitäten
  \item 25-30\,\% Thorax
  \item 10-37\,\% Abdomen und Becken
  \item 5-15\,\% Wirbelsäule
\end{itemize}

\Fazit{Die Arbeitsdiagnose Polytrauma kann auch ausschließlich aufgrund des 
Unfallmechanismus und der Vitalparameter gestellt werden. Patient gilt als 
polytraumatisiert bis zum Beweis des Gegenteils! Jedes Polytrauma ist notarztpflichtig.}

Die Sterblichkeit bei polytraumatisierten Patienten liegt zwischen 20-50 \%, ist
also hoch. Ein gezieltes, strukturiertes und zügiges Vorgehen ist daher
unabdingbar! Im Idealfall sollten zwischen Eintreffen der Rettungskräfte und
Transportbeginn nicht mehr als 30 Minuten vergehen. \cite{BoehmerTaschRett6}

\MassnahmenNG{Spezielle Maßnahmen}{Polytrauma}{}
{%
\begin{itemize}
  \item Schneller Traumacheck (\prettyref{topic:traumacheck}) gleichzeitig
  mit Sichern der Vitalparameter um kritische Verletzungen rasch zu
  erkennen!
  % TODO Diskussion: Ausführlicher Traumacheck beim Polytrauma
  \opt{STATUS-EXPERIMENTAL}{\A{\par\hrule
  DISKUSSION: (Der ausführliche Traumacheck ist beim Polytrauma
  zweitrangig und soll nur durchgeführt werden wenn dadurch der
  Transport nicht verzögert wird.)\hrule\par}}
  \item \TxMassVitBes
  \begin{itemize}
    \item \ce{O2}: 15\LMin
    \item Lagerung: situationsgerecht, entsprechend den Verletzungen
    und der Herz-Kreislauf-Situation -- Prioritäten absteigend:
    {
    \scriptsize
    \begin{enumerate}
      \item Atem- und Kreislaufstillstand $\rightarrow$ Rückenlage
      \item Bewusstlosigkeit, nicht intubiert $\rightarrow$ stabile Seitenlage 
      (bei Thorax- und/oder Extremitätentrauma auf verletzte Seite legen)
      \item Atemnot bei Thoraxtrauma $\rightarrow$ Oberkörper hoch (45\textdegree)
      \item Abdominaltrauma $\rightarrow$ Schonhaltung, Knierolle
      \item SHT $\rightarrow$ Oberkörper hoch (30\textdegree)
      \item Wirbelsäulentrauma $\rightarrow$ Flachlagerung auf Vakuummatratze
      \item Extremitätentrauma $\rightarrow$ Extremität schienen und fixieren
      \item Schock $\rightarrow$ Schocklagerung durch schräg stellen des Tragentisches
    \end{enumerate}
    }
  \end{itemize}
  \item Starke Blutungen stillen.
  \item \Es{Wärmeerhalt}: Der Wärmeerhalt hat massiven Einfluss auf das
  Überleben!
  \item wenn einigermaßen stabil, gründlicher  Traumacheck
  \item Verletzungsmuster/Unfallmechanismus erheben $\rightarrow$ zu 
  erwartende Komplikationen einschätzen
  \item \Z{Notarzt: Schmerztherapie, Narkose, Intubation und Beatmung}
  \item Transport:
  \begin{itemize}
    \item Aviso im Zielspital (über Leitstelle)
    \item bei Wirbelsäulen-Trauma Rücksprache mit der Leitstelle wegen ev. Hubschrauber-Transport
  \end{itemize}
\end{itemize}
}

% UNFALLORT IN DEN ABSCHNITT EINSATZTAKTIK VERSCHOBEN (Emhofer)

%\subsubsection{Der Unfallort und der Unfallmechanismus}

%\begin{itemize}
%    \item Absichern  der Unfallstelle,  möglichst auffällig machen (Blaulicht,
% Warnblinkanlage, Scheinwerfer)
% \item bei mehreren Verletzten: Triage
% \item Nachalarmierung : andere  Rettungsmittel, Feuerwehr, Polizei
% \item Eingeklemmte im Fahrzeug stabilisieren,  nicht eigenmächtig bergen
% \item Crash -Rettung nur  bei  Fremdgefährdung, HWS stabilisieren!
%\end{itemize}

\paragraph{Polytrauma beim Kind} Der Schockindex ist beim Kind unbrauchbar, da die
Normalwerte andere sind! Kinder bleiben lange stabil, verfallen dann aber schlagartig!


%%%%%%%%%%%%%%%%%%%%%%%%%%%%%%%%%%%%%%%%%%%%%%%%%%%%%%%%%%%
%%%%%%%%%%%%%%%%%%%%%%%%%%%%%%%%%%%%%%%%%%%%%%%%%%%%%%%%%%%
%%%%%%%%%%%%%%%%%%%%%%%%%%%%%%%%%%%%%%%%%%%%%%%%%%%%%%%%%%%

\section{Extremitätentrauma und Sportverletzungen}
%Gelenke betreffend
Verletzungen der Extremitäten können in verschiedenen Formen auftreten:
\begin{itemize}
  \item Verstauchung (Distorsion)
  \item Verrenkung (Luxation)
  \item Knochenbruch (Fraktur)
  \item Amputationen
\end{itemize}

\subsection{Verstauchung (Distorsion), Bänderzerrung und Bänderriss}
\Definition{Zerrung/Überdehnung des Kapsel-/Bandapparates}

Bei der Verstauchung (\T{Distorsion}) handelt es sich um eine Gelenksverletzung,
bei welcher ein Knochen durch eine stumpfe Gewalteinwirkung kurz verschoben oder
verdreht wird. Im Gegensatz zur Verrenkung springt der Knochen aber von selbst
wieder in seine ursprüngliche Stellung im Gelenk zurück.

\paragraph{\TxSymptome}
\begin{itemize}
  \item Schmerzen (vor allem Bewegungs- und Belastungsschmerz)
  \item Schwellung
  \item Bluterguß (Hämatom)
  \item ev. Bewegungseinschränkungen
  \item hörbares Schnalzen beim Bänderriss
\end{itemize}

\MassnahmenNG{Spezielle Maßnahmen}{Verstauchung}{}
{%
\begin{itemize}
  \item Kleidung und Schmuck von der betroffenen Gliedmaße entfernen
  ($\rightarrow$ Schwellung)
  \item ruhigstellen
  \item hochlagern
  \item kühlen
  \item ad Unfallabteilung (muss nicht sofort sein)
\end{itemize}
}





\subsection{Verrenkung (Luxation)}
\label{topic:verrenkung}

\Definition{\T{Luxation}: "`Verrenkung." Gelenkverletzung, bei der zwei das
Gelenk bildende Knochenenden aus ihrer normalen Stellung zueinander verschoben
werden. Dabei kommt es oft zu Begleitverletzungen an der Gelenkkapsel bzw. den
Gelenkbändern. \cite{Gorgass7}}

Die Verrenkung (\T{Luxation}) ist eine Gelenksverletzung, bei der -- im Gegensatz
zur Verstauchung -- der verschobene oder verdrehte Knochen nicht von selbst in
das Gelenk zurück gesprungen ist. Die betroffene Extremität kann daher nicht mehr
bewegt werden. Bei einer Luxation dürfen keine eigenständigen Einrenkversuche
unternommen werden!

\paragraph{\TxSymptome}
\begin{itemize}
  \item Schmerzen
  \item Schwellung
  \item Bluterguss
  \item abnorme Gelenkstellung (federnd fixiert), bewegungsunfähig
\end{itemize}

\MassnahmenNG{Spezielle Maßnahmen}{Verrenkung}{}
{
\begin{itemize}
  \item Kleidung und Schmuck von der betroffenen Gliedmaße entfernen 
  ($\rightarrow$ Schwellung)
  \item Stellung der Gliedmaße belassen
  \item federnd fixierte Lagerung
  \item \Es{keine Einrenkversuche!}
\end{itemize}
}


%%%%%%%%%%%%%%%%%%%%%%%%%%%%%%%%%%%%%%%%%%%%%%%%%%%%%%%%%%%
%%%%%%%%%%%%%%%%%%%%%%%%%%%%%%%%%%%%%%%%%%%%%%%%%%%%%%%%%%%
%%%%%%%%%%%%%%%%%%%%%%%%%%%%%%%%%%%%%%%%%%%%%%%%%%%%%%%%%%%

\subsection{Knochenbrüche (Frakturen)}
\label{topic:frakturen}

Einen Knochenbruch nennt man auch \T{Fraktur}. Man erkennt ihn am  häufigsten
durch eine Fehlstellung des Knochens und durch abnorme Beweglichkeit  der
betroffenen Gliedmaße. Seltener sind freie Knochenfragmente zu sehen.  Je nach
betroffenen Knochen kann es zu unterschiedlich großen Blutverlusten kommen:
\begin{labeling}[~~:]{Unterschenkel}
  \item[Oberarm] \Es{\VonBis{100}{800}\,ml}
  \item[Unterarm] \Es{\VonBis{50}{400}\,ml}
%	\item[Rippen: \Es{2000-3000\,ml
  \item[\Ca{Becken}] \Es{\VonBis{500}{\Ca{5\,000}}\,ml} (Schaftfraktur!)
  \item[\Ca{Oberschenkel}] \Es{\VonBis{300}{\Ca{2\,000}}\,ml}
  \item[Unterschenkel] \Es{\VonBis{100}{1\,000}\,ml}
\end{labeling}
Die Blutung kann aus dem Knochenmark, zerrissener Muskulatur oder aus einem 
verletzten Gefäß kommen.

\subsubsection{Arten, Symptome und allgemeine Maßnahmen bei Knochenbrüchen}

\paragraph{Arten}
\begin{itemize}
  \item geschlossene Fraktur
  \item offene Fraktur (aufgrund der Gewalteinwirkung kommt es auch zu  einer
  Wunde oberhalb der Bruchstelle)
\end{itemize}

% % gestrichen lt. STEUER 2009-08-29
% Weiter unterscheidet man:
% \begin{itemize}
%     \item komplette Fraktur (vollständige Durchtrennung)
%     \item inkomplette Fraktur (unvollständige Durchtrennung)
% \end{itemize}


\paragraph{Gefahren}
\begin{itemize}
  \item starke Blutungen, ev. bis Schock (z.\,B. Oberschenkelfraktur)
  \item begleitende Nervenverletzungen
  \item Fettembolie (Fraktur langer Röhrenknochen)
  \item Infektionsgefahr. Bei offenen Frakturen ist Keimfreiheit be\-sonders 
  wichtig, da  die Gefahr einer lebenslangen Knochen\-marks\-ent\-zündung 
  \Z{(Osteomyelitis)} besteht!
  Wenn mögllich soll eine Wundfolie (z.\,B.
  \index{OpSite}\T{OpSite}\texttrademark ) verwendet werden.    \item
  \item Kombination von mehreren Frakturen $\rightarrow$ großer Blutverlust
  \item schwere Muskelquetschung kann zu einem Sauerstoffmangel im Muskel 
  führen (Kompartmentsyndrom)
\end{itemize}

Auch bei einer harmlos anmutenden Fraktur muss ein gründlicher Traumacheck
durchgeführt werden, damit keine Begleitverletzungen übersehen werden. Dafür
ist es unbedingt notwendig, dass der Patient entkleidet wird.\footnote{Nachdem
der Patient für die Untersuchung entkleidet wurde, muss er nach der
Untersuchung zugedeckt werden (Wärmeerhaltung)!} Beim Entfernen der Schuhe ist
darauf zu achten, dass ein Sanitäter das Bein fixiert während der andere
Sanitäter mittels Stiefelgriff (siehe \prettyref{topic:vacuumbeinschiene-anwendung})
vorsichtig den Schuh öffnet und schließlich vom Fuß abzieht.

\paragraph{Frakturzeichen} \index{Frakturzeichen}
\label{topic:frakturzeichen}

Es gibt sichere und unsichere Frakturzeichen:

\begin{table}[H]
\Captionof{table}{Frakturzeichen.}

\begin{tabularx}{\linewidth}[H]{XX}\addlinespace
\multicolumn{1}{c}{\Es{Sichere Frakturzeichen}}
 &
\multicolumn{1}{c}{\Es{Unsichere Frakturzeichen}}
\\\midrule
\begin{itemize}
    \item Abnorme Beweglichkeit
    \emph{("`Pseudogelenk"')}\footnote{"`Pseudogelenk"':
    \emph{"`Wie ein Gelenk"'.} Beim Pesudogelenk hat es den
    Anschein, als würde ein Gelenk bestehen (welches nicht
    dorthin gehört).)}
    \item Fehlstellung
    (abnorme Stellung,
Verkürzung, Treppenbildung)
    \item Knochenreiben (Krepitus) \footnote{\so{NICHT}
    absichtlich ausprobieren!}
    \item Sichtbare freie Knochenteile (= offener Bruch)
\end{itemize}
 &
 \begin{itemize}
     \item Schmerz
     \item Schwellung
     \item Hämatom
     \item gestörte Funktion / Bewegungsunfähigkeit
 \end{itemize}
\\\bottomrule
\end{tabularx}
\end{table}

\begin{figure}[h]
\begin{center}
\includegraphics[width=0.4\linewidth]{./images/fraktur_fehlstellung_nativ.jpg}
\vspace{0.2\linewidth}
\includegraphics[width=0.4\linewidth]{./images/fraktur_fehlstellung_xray.jpg}
\end{center}
\Captionof{figure}{Fehlstellung -- In natura und in
Röntgendarstellung. \SourcePicture{Hauer}}
\end{figure}

\paragraph{Spezielle Maßnahmen}

Hauptaugenmerk bei der Versorgung der Knochenbrüche liegt auf der Blutstillung,
der Schienung, sowie auf dem Erkennen von großen Blutverlusten.

\MassnahmenNG{Spezielle Maßnahmen}{Fraktur}{}{
\begin{itemize}
  \item Kleidung und Schmuck von der betroffenen Gliedmaße entfernen 
  ($\rightarrow$ Schwellung)
  \item Blutstillung, Wundversorgung
  \item Bei offenen Frakturen ist Keimfreiheit be\-sonders wichtig, da  die Gefahr
  einer lebenslangen Knochen\-marks\-ent\-zündung \Z{(Osteomyelitis)}
  besteht! Wenn möglich soll eine \E{Wundfolie} (z.\,B.
  \index{OpSite}\T{OpSite}\texttrademark ) verwendet werden.
  \item Bei starken Schmerzen oder großem Blutverlust $\rightarrow$ Notarzt
  \item Ruhigstellung/Schienung
  \item Ggfs. Schockbekämpfung
\end{itemize}
}

\subsubsection{Spezielle Knochenbrüche}

\paragraph{Schlüsselbeinfraktur} Die Schlüsselbeinfraktur entsteht meistens, wenn
der Patient mit ausgestreckten Armen stürzt. Die große Gefahr dieser Fraktur
liegt darin, dass die Schlüsselbeinarterie und -vene verletzt werden können. Die
Ruhigstellung erfolgt einerseits durch ein Armtragetuch (Dreiecktuch) und
andererseits durch Fixation des Arms am Oberkörper durch ein zweites Dreiecktuch.

\paragraph{Schenkelhalsfraktur} Bei der Schenkelhalsfraktur handelt es sich um
einen Bruch des Oberschenkelhalses. Meist sind ältere Menschen betroffen, die auf
die Hüfte gestürzt sind. Symptome sind:
\begin{itemize}
	\item starker Bewegungs- oder Stauchungsschmerz
	\item betroffenes Bein ist nach außen gedreht und verkürzt
	\item Unfallhergang meist Sturz (auch an Kollaps denken!)
\end{itemize}

Die Versorgung erfolgt durch Rettung mittels Schaufeltrage und Ruhigstellung
mittels Vakuummatratze. An große Blutverluste denken (Kontrolle der
Vitalparameter)! Eventuell kann eine Schmerztherapie (Notarzt) notwendig sein.

\subsection{Amputationen}

\Definition{Amputation bezeichnet einen kompletten oder inkompletten  Abriss
bzw. Abtrennung (von Teilen) einer Extremität oder eines anderen Körperteils, 
sodass deren Durchblutung ganz oder teilweise aufgehoben ist. Den abgetrennten 
Körperteil nennt man Amputat. \cite{BoehmerTaschRett6}}

Amputationen erfolgen meist bei Motorradunfällen oder Schnittverletzungen (z.\,B.
mit einer Kreissäge). Je glatter die Wundränder sind desto größer ist die Chance
dass das Amputat wieder erfolgreich angenäht werden kann. Der Amputatstumpf wird
steril mit einem Druckverband versorgt (kein Abbinden!). An große Blutungen und
Schockgefahr denken!

\MassnahmenNG{Spezielle Maßnahmen}{Versorgung eines Amputats}{}
{%
\begin{itemize}
  \item keine Reinigung des Amputats
  \item sterile Versorgung des Amputats (Verband)
  \item Lagerung und Transport in doppelwandigem Amputatsbeutel bei  einer
  Kälte von ca. 4°C. Zwischen die beiden Schichten des Beutels gibt  man Wasser
  mit Eis (im Verhältnis 1:1).
  \item Das Amputat darf auf keinen Fall gefroren werden oder  direkten Kontakt
  zur kühlenden Substanz haben!
\end{itemize}
}

\subsection{Weitere Verletzungsbilder}

Neben den bereits genannten Verletzungen, treten (im Sport oder in der Freizeit)
oft auch folgende, meist harmlose, Verletzungen auf, welche hier der
Vollständigkeit halber erwähnt sind:

\begin{itemize}
  % 	\item Muskelkrampf
     \item Muskelzerrung
     \item Muskelfaserriss
     \item Achillessehnenriss
     \item Meniskusschaden
\end{itemize}

%%%%%%%%%%%%%%%%%%%%%%%%%%%%%%%%%%%%%%%%%%%%%%%%%%%%%%%%%%%
%%%%%%%%%%%%%%%%%%%%%%%%%%%%%%%%%%%%%%%%%%%%%%%%%%%%%%%%%%%
%%%%%%%%%%%%%%%%%%%%%%%%%%%%%%%%%%%%%%%%%%%%%%%%%%%%%%%%%%%
%%%%%%%%%%%%%%%%%%%%%%%%%%%%%%%%%%%%%%%%%%%%%%%%%%%%%%%%%%%

\section{Verbrennung}
\label{topic:verbrennung}

\Definition{Die Verbrennung ist eine durch Hitzeeinwirkung entstandene
Schädigung bzw. Verwundung der Haut.}

\marginlabel{\cite{Kamolz2009,Verbrennungsmed2009Praehosp,Beneker2004,Loennecker2001}}
Verbrennungen reichen von kleinen Brandblasen bis schweren Schädigungen der Haut
und der darunter liegenden Gewebe (z.\,B. großflächige Verkohlung). Um die Schwere
der Verletzung rasch einschätzen zu können muss die Ausdehnung und der
Schweregrad der Verbrennung beurteilt werden. Dazu stehen die Prozentregel sowie
die Definitionen der Schweregrade zur Verfügung.




\Layout{22}{Ausdehnung: Handregel und Neunerregel}
{%
Die Größe des geschädigten Areals wird in Prozent der
ver\-brannten Körperoberfläche angegeben. Diese kann entweder mit der
\T{Handregel oder mittels der
} \T{Neunerregel} ermittelt werden. Die Handregel besagt, dass die Handfläche
des Patienten einem Prozent seiner Körperoberfläche entspricht. 

Da dies in der Praxis bei großflächigen Verbrennungen meist umständlich ist, gibt
es auch noch die \Es{Neunerregel}. Hier wird der Körper in
9\,\%-Teile eingeteilt um das Ausmaß der Vebrennung abschätzen zu können.

\begin{center}
\begin{tabulary}{\linewidth}[H]{LCC}
\toprule
 & \% &
\\\midrule
Kopf:
 & 9  &
\\
Arm, je:  &  9  &
\\
Rumpf vorne: & 18  &
\\
Rumpf hinten: & 18 &
\\
Bein, je: & 18 &
\\
Genitalien:  & 1 &
\\\bottomrule
\end{tabulary}
\end{center}
% \Captionof{table}{Neunerregel} \label{tab:Neunerregel}

\TODO{GABS}{2010-04-18}{Bild: Neunerregel}{}
}
{\begin{itemize}
    \item Handregel
    \item Neunerregel
\end{itemize}}
{./images/noauth/AASdev20090228-img94.png}
% {./icons/under-construction.png}
{Neunerregel.\label{tab:Neunerregel}}{}

Die Schweregrade einer Verbrennung werden wie folgt eingeteilt:
\begin{center}
\begin{tabulary}{\linewidth}[H]{CL}
\toprule
 Grad
 & Symptome
\\\midrule
1\,\textdegree
 &
Rötung, Schmerz, Schwellung
\\
2\,\textdegree
 &
Blasenbildung, Schmerz
\\
3\,\textdegree
 &
\T{Verschorfung}, Schmerzen
abhängig von der Tiefe der Schädigung
(Zerstörung von Nerven), Schrumpf\-ung des Gewebes durch
Hitze\-ein\-wirkung, reicht bis zur Verkohlung 
\\\bottomrule
\end{tabulary}
\end{center}

Mit den erhobenen Werten (\% der verbrannten Körperoberfläche und Schweregrad)
beurteilen Sie die Schockgefahr wie folgt:

\Cave{Ab \Es{5\%} verbrannter (mind. 2. Grad) Körperoberfläche bei
Kindern oder \Es{10\%} bei Erwachsenen besteht Schockgefahr
(Flüssigkeitsverlust)!}

% TODO INHALT: Verbrennung:, leichte, Maßnahmen
% Schockbekämpfung?
% Notarzt? wann leichte Verbennung?

\MassnahmenNG{Spezielle Maßnahmen}{Verbrennung, leicht}{}
{    
\begin{itemize}
  \item Eigenschutz!
  \item Kleiderbrand löschen
  \item Kleidung rasch entfernen (sofern nicht  eingeschmolzen!)
  \item mind. 10 min. mit mäßig kaltem Wasser die betroffenen Körperregionen spülen  --
  \Es{nicht unterkühlen!}
  \item keimfreier, nicht verklebender Verband, locker anlegen
  % DISCUSS INHALT: Water-Jel, Stellungsnahme
  %         \item wenn vorhanden: Water-Jel
  \item Schockbekämpfung
  \item Notarzt wenn nötig
\end{itemize}
    
    \Cave{CAVE Unterkühlung bei zu langer oder zu kalter
    Spülung!}
}

\MassnahmenNG{Spezielle Maßnahmen}{Verbrennung, großflächig}{}
{
\begin{itemize}
  \item \TxMassVitBes
  \begin{itemize}
    \item \ce{O2}: 10\LMin
  \end{itemize}
  \item Kleidung rasch entfernen (sofern nicht  eingeschmolzen!)
  \item mind. 10 min. mit mäßig kaltem Wasser die betroffenen Körperregionen spülen  --
  \Es{NICHT unterkühlen!}
  \item keimfreier, nicht verklebender Verband, locker anlegen
  % TODO INHALT: Stellungsnahme zu Water-Jel
  %         \item wenn vorhanden: Water-Jel
  %         \footnote{Streitfrage: Water-Jel: Umstritten}
  \item Schockbekämpfung
\end{itemize}

\Cave{CAVE Unterkühlung bei zu langer oder zu kalter Spülung!}
}

\Layout{2}{Was man \Ca{nicht} tun sollte \ldots}
{%
\begin{itemize}
    \item[\Irritant] NEIN: Salbe, Puder, etc. auf frische  Brandwunde
    \item[\Irritant] NEIN: Blasen öffnen
    \item[\Irritant] NEIN: festklebende Kleidung gewaltsam  entfernen
    \item[\Irritant] NEIN: aluminiumbeschichteten Verband  verwenden (Gefahr der
    Hitzereflexion bei  nicht ausreichend gekühlten  Verbrennungen!)
\end{itemize}
}{

}{}{}{}







\subsection{Inhalationstrauma}
\label{topic:inhalationstrauma}
\index{Inhalationstrauma}

\Definition{Von einem Inhalationstrauma spricht man bei
Einatmung von reizenden, giftigen und/oder heißen Gasen,
welche eine Schädigung der Atemwege und der Lunge zur Folge
hat.}

\paragraph{Leitsymptome}

\begin{itemize}
  \item Z.\,n. Gasexposition
  \item angebrannte Gesichtshaare
  \item geschwollene, rußige Lippen
  \item Rußpartikel im Auswurf
  \item Husten
  \item Heiserkeit
  \item Stridor
  \item Atemnot
\end{itemize}

Das volle Ausmaß der Schädigung kann unmittelbar oder
innerhalb von Stunden erreicht werden. Eine engmaschige
Beobachtung des Patienten ist daher sehr wichtig.

\MassnahmenNG{Spezielle Maßnahmen}{Inhalationstrauma}{}{

\begin{itemize}
  \item Vitale Bedrohung einschätzen. \par
  \TxBeiVit \TxMassVitBes
  \begin{itemize}
    \item \ce{O2}: 15\LMin
  \end{itemize}
  \item Lagerung: Oberkörper hoch
  \item \ce{O2}: \VonBis{6}{8}\LMin
  \item Engmaschiges Monitoring
\end{itemize}

\ASBFLO{\Z{Notfallsanitäter geben bei einer Rauch- oder
Reizgasinhalation ein inhalatives Kortikoid (z.\,B.
Pulmicort\texttrademark) laut Algorithmus der
Medikamentenliste 1.}}

}









%%%%%%%%%%%%%%%%%%%%%%%%%%%%%%%%%%%%%%%%%%%%%%%%%%%%%%%%%%%
%%%%%%%%%%%%%%%%%%%%%%%%%%%%%%%%%%%%%%%%%%%%%%%%%%%%%%%%%%%

\subsection{Verletzungen durch chemische Stoffe -- "`Chemische Verbrennungen"'}
\label{topic:veraetzung-trauma}

\noindent\dsliterary{} \cite{RD200811Hoffmann}

Vergiftungen durch Einnahme von Säuren und Laugen werden 
unter \prettyref{topic:veraetzung-intoxikation} besprochen. 

\begin{wrapfigure}{r}{4cm}
\includegraphics[width=4cm]{./images/trauma/notbrause2.jpg}
\end{wrapfigure}
Schädigungen der Haut durch Säuren, Laugen oder andere schädliche Stoffe weisen
-- trotz der unterschiedlichen Entstehung -- Gemeinsamkeiten mit
Verbrennungsschäden auf. Daher werden diese Verletzungen oft auch als
\emph{"`chemische Verbrennung"'} bezeichnet, auch wenn meist keine Hitze im
Spiel ist.  

Genau wie bei der Verbrennung kommt es zu einem Ausfall der Schutzfunktion der
Haut, oft zu Schmerzen und zu gefährlichen Schockzuständen. Auch die
Einschätzung der Ausdehnung der Verwundung erfolgt wie bei der Verbrennung mit
der Neuner-Regel o.\,ä. 

Eine Besonderheit stellt jedoch der verantwortliche Stoff dar, mit welchem u.U.
der Patient, dessen Kleidung und evtl. auch die Umgebung kontaminiert ist. Von
diesem Stoff kann sowohl für den Patienten, als auch für den Helfer eine Gefahr
ausgehen. Auf den Eigenschutz und die Dekontamination muss großer Wert gelegt
werden. 

\MassnahmenNG{Spezielle Maßnahmen}{Verletzungen mit chemischen
Substanzen}{}{

\begin{enumerate}
  \item Auf Selbsschutz achten!
  \begin{itemize}
    \item Patienten ggfs. retten, aber nur wenn es der
    Eigenschutz zulässt.
  \end{itemize}
  \item Gefahr erkennen: \Es{Welcher Stoff?}  Expertenrat einholen!
  \item Wenn erforderlich: Absperrmaßnahmen durchführen
  \item Wenn erforderlich:  Spezialkräfte anfordern.
  \item Vorgehen analog zu Abschnitt "`Gefahrengutunfall"'
  (\prettyref{topic:gefahrengutunfall})
  \item Weitere Opfer ausschließen (lassen)
  \item Sicherung der Vitalfunktionen, Beurteilung ob der
  Patient vital bedroht ist.
  \item \TxBeiVit \TxMassVitBes
  \begin{itemize}
    \item Lagerung: situationsgerecht.
  \end{itemize}
  \item Entfernung kontaminierter Kleidung
  \item Spülung der betroffenen Hautpartien mit Wasser  (richtig abrinnen
  lassen)
  \item Auge: Spülung mit Wasser / Ringerlösung (nach außen abrinnen
  lassen!),  Kontaktlinsen entfernen
  \item Hospitalisation: Ab einer Ausdehnung von 15\,\% (Erwachsener) bzw.
  10\,\% (Kind) Spezialabteilung für Verbrennungsverletzungen. Bei Augenverletzungen
  Abt. für Augenheilkunde.
\end{enumerate}
}










%%%%%%%%%%%%%%%%%%%%%%%%%%%%%%%%%%%%%%%%%%%%%%%%%%%%%%%%%%%
%%%%%%%%%%%%%%%%%%%%%%%%%%%%%%%%%%%%%%%%%%%%%%%%%%%%%%%%%%%
%%%%%%%%%%%%%%%%%%%%%%%%%%%%%%%%%%%%%%%%%%%%%%%%%%%%%%%%%%%

\section{Erfrierungen}
\label{topic:erfrierung}

\Definition{...sind Gewebsnekrosen, die durch lokale 
Durchblutungsstörungen aufgrund von  Kälteeinwirkung (Eis, Wind, Nässe,  Kälte)
entstanden sind.}

Besonders gefährdet sind körperferne Körperteile wie Zehen,
Finger, Ohren oder Nase.

Einteilung erfolgt in drei \Es{Stadien} (siehe 
\prettyref{tab:Erfrierungen}), diese haben am Notfallort keine entscheidende
Bedeutung für die zu setzenden Maßnahmen. Den Schweregrad einer Erfrierung kann man erst nach Tagen voll
einschätzen. Deshalb ist es wichtig, die Warnsymptome ernst zu nehmen!

\begin{center}
\begin{tabulary}{\linewidth}[H]{CLC}
\toprule
 Stadien
 & Symptome
 &
\\\midrule
1
 &
Blässe, Empfindungsstörungen
 &
\\
2
 &
Blasenbildung, Schmerzen
 &
\\
3
 &
grau/schwarze Marmorierung (Nekrose), Gefühllosigkeit
\\\bottomrule
\end{tabulary}
\Captionof{table}{Stadien der Erfrierung \label{tab:Erfrierungen}} 
\end{center}


\MassnahmenNG{Spezielle Maßnahmen}{Erfrierungen}{}
{
\begin{itemize}
  \item beengende Kleidungsstücke lockern (Einschnürung erhöht
  Risiko einer Erfrierung)
  \item Keimfreier, gepolsterter Verband
  \item Warme, gezuckerte Getränke -- KEIN Alkohol
  \item Körperstamm wärmen, NIE das erfrorene Körperteil abreiben!!
\end{itemize}
}




%%%%%%%%%%%%%%%%%%%%%%%%%%%%%%%%%%%%%%%%%%%%%%%%%%%%%%%%%%%
%%%%%%%%%%%%%%%%%%%%%%%%%%%%%%%%%%%%%%%%%%%%%%%%%%%%%%%%%%%
%%%%%%%%%%%%%%%%%%%%%%%%%%%%%%%%%%%%%%%%%%%%%%%%%%%%%%%%%%%

\section{Unfälle durch Strom- und Blitzschlag}

\marginpar{\includegraphics[width=\linewidth]{./images/electricity_japan-00800.jpg}
\small\emph{Foto: Fina}}



Stromunfälle kommen typischerweise in folgenden Situationen vor:
\begin{itemize}
  \item Spielende Jugendliche oder Arbeiter auf Eisenbahnwaggons
  \item LKW-Hebekräne oder Bagger geraten in Freileitungen
  \item Unachtsamer Umgang mit Elektrogeräten im Badezimmer
  \item Unsachgemäße Reparatur von Elektrogeräten
  \item Arbeitsunfälle (Elektriker, Heimwerker)
  \item Tragen von  Metallteilen in der Nähe von Hochspannungsleitungen
\end{itemize}


\Layout{2}{Allgemeines zum elektrischen Strom und seiner
Wirkung auf den Körper}
{%
Die Gefährdung bei Stromunfällen hängt von verschiedenen Faktoren ab:
\begin{description}
  \item[Elektrische Spannung:] Die elektrische Spannung
  ist die Ursache für den elektrischen Stromfluss und wird
  in der Einheit \T{Volt} (V) gemessen. Wie gefährlich eine Spannung ist hängt
  nicht nur von der Höhe der Spannung selbst, sondern wesentlich auch vom 
  elektrischen Strom, der von ihr hervorgerufen wird, ab!\footnote{So kann z.\,B.
  die 12\,V-Spannung einer Modelleisenbahn völlig ungefährlich sein, gibt es
  jedoch im Auto bei einer (großen) 12\,V-Batterie einen Kurzschluss, kann ein
  gefährlich großer Strom fließen.}
  \begin{itemize}
    \item \T{Niederspannung}: Von Niederspannung spricht man bei Spannungen
    bis 1000\,V. \ZFootnote{An einer Autobatterie liegt eine Spannung von 12 V. An
    einer Steckdose liegt eine Spannung von 230 V, in Industriebetrieben wird
    zusätzlich oft eine Spannung von 400 V verwendet. Diese Spannungen sind auf
    jeden Fall lebensgefährlich. Ein Defibrillator schockt mit etwa 1000\,V
    (!), ein externer Herzschrittmacher arbeitet mit bis zu 400\,V.}
    \item \T{Hochspannung} \ldots\ ist eine Spannung über 1000\,V.
    \ZFootnote{Über Hochspannungsleitungen wird die elektrische Energie mit
    einer Spannung von 110\,000\,V (= 110\,kV) bis 380\,000\,V (= 380\,kV) übertragen.
    An den Oberleitungen der Eisenbahn liegen 15\,000\,V (= 15\,kV). Bei
    Gewitterblitzen treten Spannungen von 10-30 Millionen Volt (10-30 Megavolt) auf.}
  \end{itemize}
  \item[Stromstärke:] Der elektrische 
  Strom\ZFootnote{Elektrischer Strom ist der Fluss von Elektronen.  Die
  treibende Kraft hinter dem Stromfluss ist die elektrische Spannung.}  wird
  durch die elektrische \T{Stromstärke} beschrieben, welche mit der  Einheit
  \T{Ampere} (A) gemessen wird. Kommt es  zu einem Stromunfall ist die 
  Schädigung im Körper von der Höhe der Stromstärke  abhängig.\ZFootnote{Bei
  geringen Stromstärken \E{unter 1 Milliampere (mA)}  ist  nur ein
  elektrischer Schlag spürbar. Bei Stromstärken \E{ab ca.  10 Milliampere
  (mA)} kommt es zu Schmerzempfindungen und  Muskelkontraktionen \cite{Jaros1},
  ab \Es{15-25 mA} kommt es aufgrund der Muskelkontraktionen zu  krampfhaften
  Festhalten der berührten Stromquelle \cite{BoehmerTaschRett6}.  Stromstärken ab
  \E{25-50\,mA} sind \E{lebensgefährlich}!~\cite{BoehmerTaschRett6}
  Ab \E{80\,mA} kommt es zu Atemlähmungen und
  Kammerflimmern \cite{Jaros1}, Stärken über \E{3\,A} führen zusätzlich zu
  Verbrennungen, Verkochungen und Verkohlungen. \cite{LPN3} Die Werte beziehen 
  sich auf Wechselstrom, die Grenzwerte für Gleichstrom sind höher. Da oft die 
  Stromart nicht bekannt ist, ist es prinzipiell sinnvoll von der
  gefährlicheren  Stromart auszugehen!} Stromstärken ab 25\,mA gelten als
  potentiell tötlich!
  \item[Widerstand:] Jeder Körper und
  jeder Stoff setzt dem elektrischen Stromfluss einen
  Widerstand entgegen. Dieser elektrische Widerstand wird in
  \T{Ohm} ($\Omega$) gemessen. Je größer der Widerstand eines
  Körpers, desto kleiner ist der elektrische Strom. Trockene
  Haut hat einen elektrischen Widerstand von 100\,000 Ohm,
  \Es{feuchte Haut hat einen geringeren Widerstand} von
  weniger als 1\,000 Ohm. Der Widerstand von geschädigter
  (z.\,B. verbrannter) Haut ist noch geringer.\ZFootnote{Greift man z.\,B. mit
  feuchten Händen in die Steckdose (230\,V) so stellt sich durch den 
  Körperwiderstand von ca. 1000 Ohm ein elektrischer Strom der Stärke 23\,mA 
  ein. Diese Stromstärke ist absolut  lebensgefährlich, wie man aus der
  Übersicht  oben schnell entnehmen  kann.}\ZFootnote{Man kann dies  durch 
  das Ohmsche Gesetz leicht selbst nachrechnen! Die Spannung $U$ ist das 
  Produkt aus dem Widerstand $R$ und der elektrischen Stromstärke $I$: $U=R \cdot I$}
  \item[Einwirkdauer:] Je länger die Einwirkzeit (Kontaktzeit), desto
  größer ist die Schädigung.
  \item[Stromweg:] Die Schädigung im Körper hängt auch vom Weg des Stromes
  durch den Körper ab. Fließt der Strom z.\,B. von einer Hand zur anderen Hand
  durch den Körper ist dies gefährlicher, als wenn der Strom von einer Hand zum
  Fuß fließt.
  \item[Stromart:] Wechselstrom ist gefährlicher als Gleichstrom. An den
  Steckdosen im Haushalt liegt Wechselspannung an. Dieser ändert die
  Stromrichtung 50~mal pro Sekunde (\Z{Frequenz = 50\,Hz}), weshalb die
  Wahrscheinlichkeit für Kammerflimmern auch bei kurzen Stromkontakten steigt.
  \cite{RedelsteinerHRNS1} \ZFootnote{Bei monophasischen Defibrillatoren wird
  mit Gleichstrom gearbeitet, biphasische Defibrillatoren polen die Stromrichtung
  einmal um (Wechselstrom).}
\end{description}

\AuthNotes{\#9: EMHOFER: Bilder von verschiedenen
Steckdosen, Kabel und Hinweisschildern einfügen}

%\T{Niederspannung}: bis 1000\,Volt (Steckdose:  230 Volt)
%\T{Hochspannung}: über 1000\,Volt, Stromstärke
%{\textgreater}  3 Ampere (Überlandleitungen) 

}{

}{}{}{}


\Layout{2}{Besonderer Gefahrenbereich}
{%
\index{Gefahrenbereich!Strom}\label{topic:gefahrenbereich-strom}
%
Eine Starkstrom- bzw. Hochspannungsanlage muss durch einen \Es{Fachmann}
(E-Werk, ÖBB, Betriebsleitung, ...) abgeschalten bzw. geerdet werden. 

Die Anlage muss
gegen Wiedereinschalten gesichert werden! Patient mit isolierendem
Material aus dem Stromkreis befreien. Ein Sicherheitsabstand zu
stromführenden Leitungen muss eingehalten werden! 



\subparagraph{Freileitungen}

Bei Freileitungen mit
220\,000 Volt (220\,kV) sind mindestens 4\,m Abstand zu halten.
\cite{Jaros1} 

Bei herabhängenden Hochspannungsleitungen oder
nach Blitzeinschlägen auf die Schrittspannung achten: Die
Spannung fällt im Boden nicht sofort auf 0\,Volt ab, innerhalb
eines gewissen Umkreises kann daher zwischen den beiden
Beinen am Boden eine so genannte Schrittspannung auftreten.
Befinden Sie sich innerhalb eines solchen Gebietes lassen Sie
Ihre Beine eng beisammen und springen aus diesem Bereich
heraus. Je nach Bodenbeschaffenheit können Schrittspannungen
bis zu einer Distanz von 20\,m von der herunterhängenden
Hochspannungsleitung auftreten. \cite{RedelsteinerHRNS1} 

\subparagraph{Bahnanlagen}
Bei
Notfällen auf Bahngleisen ist auf jeden Fall der
Streckenabschnitt von der ÖBB sperren zu lassen (über die
Leitstelle)! An den Masten der Oberleitungen der ÖBB sind
entsprechende Mastnummern angebracht. Diese müssen bei der Leitstelle gemeldet
werden. 

Ähnliches gilt für Stadt- und U-Bahn-Anlagen. Da die Stromversorgung oft nicht
durch eine Oberleitung, sondern durch andere Vorrichtungen in
Bodennähe (Stromschiene, etc.) erfolgt, besteht auch bei intakten Anlagen die
Gefahr eines Stromschlags, sofern die Gleisanlagen betreten werden. Als
Traktionsspannung werden oft Spannungen um die 700\,V eingesetzt. \cite{U61989}
}{

}{}{}{}


\Layout{2}{\TxSymptome}
{ Die Verletzungsmuster und Symptome variieren je nach Spannung (Volt) und
Stromstärke (Ampere):

\begin{itemize}
  \item Unfallmechanismus und Unfallhergang
  \item Reizeffekte (Muskel-/ Nervenerregung):
  \begin{itemize}
    \item Verkrampfung ({\quotedblbase}kann den stromführenden Teil nicht
    loslassen{\textquotedblleft})
    \item Sensibilitätsstörungen
    \item ZNS ${\rightarrow}$ Bewußtseinsstörung
    \item Herzrhythmusstörungen (Kammerflimmern, ev.  als Spätfolge!)
  \end{itemize}
  \item Bewusstseinsstörungen bis Krampfanfälle oder Bewusstlosigkeit
  \item Lähmungen (vorübergehend)
  \item Wärmeentstehung
  \begin{itemize}
    \item Strommarken (Ein-/ Austrittsverbrennung)
    \item Lichtbogen: oberflächliche und tiefe  Verbrennungen
    \item Blitzfiguren an der Haut
  \end{itemize}
\end{itemize}
%\begin{itemize}
%    \item Gleichstrom
%    \item Wechselstrom ${\rightarrow}$ eher  Herzrhythmusstörungen
%    \item Patient  kann Spannungsquelle ev. nicht  loslassen (Krampf der
%    Handmuskulatur)
%    \item Auch Defi stellt Strom-Gefahrenquelle  dar!!!
%    \item Blitzunfall
%    \begin{itemize}
%        \item mehrere Millionen (!) Volt
%        \item bis 100.000 Ampere
%    \end{itemize}
%\end{itemize}

%\paragraph{\TxSymptome}
%\begin{itemize}
%    \item Bewußtlosigkeit
%    \item Lähmungen (vorübergehend)
%    \item Blitzfiguren auf der Haut
%    \item Herzrhythmusstörungen möglich
%\end{itemize}

\Cave{Das eigentliche Ausmaß der Verletzung ist oft
{\quotedblbase}unsichtbar{\textquotedblleft}, da der Strom im
Körperinneren -- also {\quotedblbase}unterirdisch{\textquotedblleft}
-- viel Schaden anrichtet. }



\subparagraph{Blitzschlag} Bei Gewitterblitzen treten Spannungen von Millionen
von Volt und Stromstärken von bis zu 100\,000 Ampere auf. Trotzdem gibt es auch
bei Blitzunfällen immer wieder Überlebende, da die Einwirkdauer extrem kurz
(wenige Millisekunden) ist. Typische Symptome sind

\begin{itemize}
  \item Bewusstlosigkeit
  \item vorübergehende Lähmung,
  \item Blitzfiguren auf der Haut.
  \item Herzrhythmusstörugnen sind möglich
\end{itemize}
}{

}{}{}{}


\MassnahmenNG{Spezielle Maßnahmen}{Stromunfälle}{}{
\begin{itemize}
  \item \Es{Eigenschutz!} Bevor  Sie den Patienten berühren,
  vergewissern Sie sich, dass der Patient keinen Kontakt  mehr mit dem
  Stromkreis hat! \Es{Schalten Sie auf jeden Fall vor jeder Berührung
  den Strom ab!}
  \begin{itemize}
    \item Gefahrenbreich beachten:
    \prettyref{topic:gefahrenbereich-strom}
    \item Sicherheitsabstand
    (\prettyref{topic:gefahrenbereich-strom})
    \item Schrittspannung und Spannungskegel
    beachten \prettyref{topic:gefahrenbereich-strom}
    \item Bahnanlagen sperren/sichern lassen
    (\prettyref{topic:gefahrenbereich-strom})
    \item Schalter ausschalten, (Haupt-)Sicherung
    auslösen, Starkstrom/Hochspannung durch Fachmann  abschalten bzw.
    erden lassen
    \item Strom gegen Wiedereinschalten sichern!
    \item Patient mit isolierendem Material aus dem Stromkreis befreien
    
  \end{itemize}
  \item Beurteilung der vitalen Bedrohung.
  \par \TxBeiVit
  \TxMassVit
  \item \Es{Monitoring \&
  Reanimationsbereitschaft}: auf evtl.
  Herzrhythmusstörungen vorbereitet sein
  (bis zu 24\,h Latenz)
  \item Genaue Patientenuntersuchung ${\rightarrow}$
  Eintritts- und Austrittsmarke
  \item Hospitalisierung auch bei symptomlosen
  Patienten zur Überwachung
  \item Stromverbrennungen wie Verbrennungen versorgen
\end{itemize}
}

\section*{Weiterführende Literatur}
\noindent\SymbolLiteratur{} 
\cite{PHTLS6,
emergency9,
OHCSP7,
Helfen2008,
DRAn3,
Patzelt2005,
Ficklscherer2005,
Sachsenweger2003,
SiewertChirurgie8,
ThierbachInterhospitaltransfer1,
RueckerBildatlas1,
ArbeitstechnikenAZ1,
ScholzNotfallmedizin2}
% \nocite{PHTLS6}
% \nocite{emergency9}
% \nocite{OHCSP7}
% \nocite{Helfen2008}
% \nocite{DRAn3}
% \nocite{Patzelt2005}
% \nocite{Ficklscherer2005}
% \nocite{Sachsenweger2003}
% \nocite{SiewertChirurgie8}
% \nocite{ThierbachInterhospitaltransfer1}
% \nocite{RueckerBildatlas1}
% \nocite{ArbeitstechnikenAZ1}
% \nocite{ScholzNotfallmedizin2}



\clearpage
\ChapterBibliography
